% concatenated from subgauss_arxiv.tex
\documentclass{article}

\usepackage[utf8]{inputenc}
\usepackage[hscale=.7, vscale=.8]{geometry}
\usepackage{bbm, amsmath, amssymb, amsthm, bm}
%\usepackage{mlmodern}
\usepackage{natbib}
\usepackage[colorlinks=true, citecolor={blue}]{hyperref}
\usepackage{mathtools}

\usepackage[utf8]{inputenc} % allow utf-8 input
\usepackage[T1]{fontenc}    % use 8-bit T1 fonts
\usepackage{hyperref}       % hyperlinks
\usepackage{url}            % simple URL typesetting
\usepackage{booktabs}       % professional-quality tables
\usepackage{amsfonts}       % blackboard math symbols
\usepackage{nicefrac}       % compact symbols for 1/2, etc.
\usepackage{microtype}      % microtypography
\usepackage{xcolor}         % colors
\usepackage{pgf, pgfplots}
\usepackage{latexsym}
\usepackage{amssymb}
\usepackage{amsthm, mathrsfs, bbm}
\usepackage{amscd}
\usepackage{amsmath}
\usepackage{amsthm}
\usepackage{titlesec}
\usepackage{url}
\usepackage{blindtext}
\usepackage{svg}
\usepackage{wrapfig}
\usepackage{wrapstuff}
\usepackage{enumitem}
\usepackage{authblk}

\theoremstyle{definition}
\newtheorem*{theorem*}{Theorem}
\newtheorem{theorem}{Theorem}
\newtheorem{conjecture}{Conjecture}
\newtheorem{remark}{Remark}
\newtheorem*{lemma*}{Lemma}
\newtheorem{lemma}{Lemma}
\newtheorem*{proposition*}{Proposition}
\newtheorem{example}{Example}
\newtheorem{proposition}{Proposition}
\newtheorem{definition}{Definition}
\newtheorem{corollary}{Corollary}
\renewenvironment{rm}{\par\noindent\upshape}{\par}
\newcommand*\D{\partial}
\newcommand{\wip}{{\color{red} \textbf{in progress}}}


\newcommand*\dd[2]{\frac{d#1}{d#2}}
\newcommand*\pd[2]{\frac{\partial#1}{\partial#2}}
\newcommand*\eps{\epsilon}
\newcommand*\veps{\varepsilon}
  \DeclareMathOperator*{\argmin}{argmin}
\renewcommand*\epsilon{\varepsilon}
\renewcommand*\phi{\varphi}
\newcommand{\R}{\mathbb{R}}
\newcommand{\C}{\mathbb{C}}
\newcommand{\Z}{\mathbb{Z}}
\newcommand{\N}{\mathbb{N}}
\newcommand{\Q}{\mathbb{Q}}
\newcommand{\E}{\mathbb{E}}
\newcommand{\1}{\mathbbm{1}}
\renewcommand{\P}{\mathbb{P}}
\newcommand*\brac[1]{\langle#1\rangle}
\newcommand*\p[1]{\left(#1\right)}

\DeclareMathOperator*{\argmax}{arg\,max}

\title{Estimation of the sub-Gaussian parameter}
\author[1]{Jason Liu}
\author[1]{Min Xu\thanks{Corresponding author: \texttt{mx76@stat.rutgers.edu}}}
\author[2]{Jinchuan Xing}

\affil[1]{Department of Statistics, Rutgers University}
\affil[2]{Department of Genetics, Rutgers University}
\date{\today}
\begin{document}
\maketitle


\begin{abstract}

The sub-Gaussian parameter (also called the variance proxy) of a mean-zero random variable $X$ is defined as $\xi^2_* = \sup_{\lambda \in \mathbb{R}} L(\lambda)$ where $L(\lambda) = \frac{2}{\lambda^2} \log \mathbb{E} e^{\lambda X}$ is a weighted cumulant generating function. Despite the ubiquity of sub-Gaussian random variables, the estimation of $\xi^2_*$ has received little attention and is not yet well understood. In this work, we study a natural estimator of $\xi^2_*$ based on constrained maximization of the empirical analogue of $L$. We prove that the estimator is consistent bound the rates of convergence under assumptions on $L$: if $L$ has an maximizer, then our bound is $O_p(n^{-1/2 + \varepsilon})$ for any $\varepsilon > 0$; if the argmax of $L$ is also bounded, then the bound improves to $O_p(n^{-1/2})$. We show that our assumptions on $L$ are necessary by proving that the minimax risk over all sub-Gaussian distributions is $\Omega(1)$; imposing increasingly strong assumptions on the tail growth of $L$ yields a continuum of classes whose minimax lower bound interpolates between $\Omega(1/\log n)$ and $\Omega(1)$. Root-n rate is possible if we restrict to a subclass of distributions where $L$ attains its supremum in a bounded region, in which case our estimator is minimax optimal. If the underlying distribution is not sub-Gaussian, we show that our estimator goes to infinity with a divergence rate controlled by the tail of the distribution. Finally, we apply our estimator in a Gene Ontology (GO) enrichment study to construct p-values for a large-scale permutation test, showing that it can serve as a reliable alternative to the peaks-over-threshold approach, particularly in regimes where the peaks-over-threshold method is of uncertain validity.

 %\vspace{0.1in}

  %Despite the ubiquity of sub-Gaussian random variables in modern machine learning and statistics, the problem of estimating the sub-Gaussian parameter $\xi_*^2$ has received little attention. Here, we propose a natural estimator $\hat \xi_n^2$ for $\xi_*^2$ based on the maximization of a weighted cumulant generating function, and study its behavior in three regimes. First, if the underlying distribution is sub-Gaussian, we prove that the estimator is consistent. If the population maximizer exists, we show that it converges at rate $O_p(n^{-1/2 + \eps})$ for all $\eps > 0$, and if the population maximizers are bounded, we give a parametric rate of $O_p(n^{-1/2})$. Under stronger assumptions, we also show that $\hat \xi_n$ is asymptotically normal, and then give asymptotic confidence intervals for $\xi_*^2$. Second, we show that even for compactly supported distributions, no estimator can be uniformly consistent. If the distributions are uniformly compactly supported, the rate improves to $\Omega(1/\log n)$, and if the variance is bounded away from 0, the rate becomes $\Omega(n^{-1/2})$, which $\hat \xi^2_n$ achieves adaptively. Third, if the distribution is not sub-Gaussian, then $\hat \xi_n^2 \xrightarrow{\mathrm{a.s.}} \infty$ with divergence rate controlled by the tail of $X$. We close with an application to Ontology (GO) enrichment analysis, and show it can achieve comparable or superior results to the peaks-over-threshold approach.
\end{abstract}



\section{Introduction}

A random variable $X$ is called sub-Gaussian if there exists $\xi^2 > 0$ such that $\mathbb{E}[e^{\lambda X}] \leq e^{\lambda^2 \xi^2 / 2}$ for all $\lambda \in \mathbb{R}$; we let $\xi^2_* \geq 0$ denote the smallest such number and refer to it as the sub-Gaussian parameter of $X$, also called the variance proxy. Sub-Gaussian random variables play a central role in machine learning and statistics because they admit an exponential tail bound $\mathbb{P}(X \geq t) \leq \exp\{ - \frac{t^2}{2 \xi_*^2}\}$ (the sub-Gaussian condition implies that $\mathbb{E}X = 0$). This is in fact an equivalent characterization: any random variable satisfying such a tail bound is sub-Gaussian. Given independent and identically distributed random variables $X_1, \ldots, X_n$ from a sub-Gaussian distribution $P$, we study the estimation of the sub-Gaussian parameter $\xi^2_*$. 

Our motivation comes from using empirical sub-Gaussian tail bounds in large-scale permutation tests. Suppose we are testing $K$ hypotheses and, for each hypothesis $k$, use $M$ random permutations to generate the null test statistic samples $T^{(k)}_1, \ldots, T^{(k)}_M$. The empirical p-value is then bound below by $\frac{1}{M+1}$, so after a multiple testing correction such as Bonferroni, the power becomes trivial when the number of hypotheses is large compared to $M$; this is often the case when the permuted test statistics are expensive to compute so that $M$ cannot be too large. However, if we know a priori that the null distribution is sub-Gaussian (which holds for example if the test statistic is bounded), then we can instead plug in an estimate of $\xi^2_*$ into the sub-Gaussian tail bound to obtain p-value much smaller than $\frac{1}{M+1}$ and recover power in regimes where the empirical p-values are powerless. 

Beyond empirical tail bounds, estimating $\xi^2_*$ is also useful for testing the sub-Gaussianity of a random variable, which is of interest in cases where the theoretical risk bounds assume sub-Gaussianity or where the estimation algorithms are designed for settings where the data is sub-Gaussian. 

Our estimation approach is based on the observation that any $\xi^2 > 0$ that satisfies $\mathbb{E}[e^{\lambda X}] \leq e^{\lambda^2 \xi^2 / 2}$ must be larger than $L(\lambda) = \frac{2}{\lambda^2} \log \mathbb{E} e^{\lambda X}$ for any $\lambda \ne 0$. The function $L$ is continuous with a removable singularity at 0 so that we can define $L(0) = \text{Var}(X)$ by continuity. Therefore, $\xi^2_*$ can be equivalently defined as 

\[
\xi_*^2 := \sup_{\lambda \in \mathbb{R}} L(\lambda). 
\]

Given independent and identically distributed observations $X_1, \ldots, X_n$, we define the empirical analogue of $L$ as $L_n(\lambda) = \frac{2}{\lambda^2} \log \bigl\{ \frac{1}{n} \sum_{i=1}^n e^{ \lambda (X_i - \bar{X})} \bigr\}$ for $\lambda \neq 0$ and $L_n(0) = \frac{1}{n}\sum_{i=1}^n (X_i - \bar{X})^2$. One idea is to estimate $\xi^2_*$ by taking $\sup_{\lambda \in \mathbb{R}} L_n(\lambda)$, i.e. the sub-Gaussian parameter of the empirical distribution, but this can incur excessive estimation error due to the unstable behavior of $L_n(\lambda)$ when $\lambda$ is relatively large, e.g. of order $\sqrt{\log n}$. We thus define our estimator as $\hat{\xi}^2_n = \sup_{|\lambda| \leq C_n} L_n(\lambda)$ for a slowly diverging $C_n$ which can be taken as $(\log n)^{1/4}$. The truncation introduces a bias whose magnitude is measured by the truncation gap function

\[
\delta(C) := \sup_{|\lambda| \geq C} L(\lambda) - \sup_{|\lambda| \leq C} L(\lambda), \, \text{ for $C \geq 0$.}
\]


The function $\delta(C)$ is decreasing and has a non-positive limit as $C \rightarrow \infty$; the bias of $\hat{\xi}^2_n$ vanishes quickly if $\delta(C)$ decays at a fast rate. Although the truncation gap function is motivated by our estimator, we prove that it in fact plays a role in determining the fundamental difficulty of estimating $\xi_*^2$ through a minimax lower bound analysis. We summarize our key results as follows:
\begin{enumerate}[leftmargin=4mm]
\item \textbf{Section 2:} We prove that $\hat{\xi}^2_n$ is always consistent when the underlying distribution is sub-Gaussian. If there is $C_0 > 0$ such that $\delta(C_0) \leq 0$, then we show $|\hat{\xi}^2_n - \xi^2_*| = O_p(n^{-1/2 + \eps})$ for any $\eps > 0$. If $C_0$ satisfies $\delta(C_0) < 0$, then we further show that $|\hat{\xi}^2_n - \xi^2_*| = O_p(n^{-1/2})$. If in addition $\argmax L(\lambda) = \{\lambda^*\}$ and $L''(\lambda^*) < 0$, then we prove that $\hat{\xi}^2_n$ is asymptotically normal and construct confidence interval for $\xi^2_*$. 
\item \textbf{Section 3:} We prove that the minimax risk of estimating $\xi^2_*$ is governed by the decay rate of $\delta_P$. For $\gamma \in [0, 1/2]$, define the decreasing function
$r_\gamma(t) = 2^{\frac{1-2\gamma}{1-\gamma}}t^{\frac{2\gamma-1}{1-\gamma}}$.
The minimax risk over subclasses of distributions satisfying $\delta_P(C) \leq r_\gamma(C)$ 
is then lower bounded by $\Omega\bigl(\frac{1}{\log^{1-2\gamma}(n)}\bigr)$. At $\gamma = 1/2$, where $r_{1/2}(t) \equiv 1$ so that no assumptions are made on $\delta_P$ and the subclass coincide with the class of all sub-Gaussian distributions, the minimax risk is lower bounded by $\Omega(1)$, i.e. uniform consistency is impossible. At $\gamma = 0$, where 
$r_0(t) = 2t^{-1}$ and where the resulting subclass encompasses all distributions supported on $[-1, 1]$, the minimax lower bound improves to $\Omega(1/\log n)$. If we instead assume that there exists $C_0, \delta_0 > 0$ such that $\delta(C_0) \leq - \delta_0$, then the minimax risk is of order $n^{-1/2}$ and our estimator $\hat{\xi}^2_n$ is minimax optimal. 
\item \textbf{Section 4:} We prove that $\hat \xi_n^2$ diverges almost surely when the underlying distribution is not sub-Gaussian, where the rate of divergence is faster if the departure from sub-Gaussianity is more significant, i.e. if the distribution is more heavy-tailed.  
\end{enumerate}

\noindent \textbf{Related Work}: The most relevant work is \cite{mies2026empiricalorlicznorms}, which considers the estimation of the Orlicz norm 
\[
\sigma_\psi \equiv \|X\|_{\psi} := \inf \{\sigma > 0: \E \psi(|X|/\sigma) < 1\}
\]
where $\psi: [0, \infty) \to [0, \infty)$ is any increasing convex function with $\psi(0) = 0$ and $\lim_{x \to \infty} \psi(x) = \infty$. 
They study the estimator $\hat \sigma_\psi := \inf \{\sigma > 0: \frac1n\sum_{i=1}^n\psi(|X_i|/\sigma) < 1\}$, i.e. the Orlicz norm of the empirical distribution, and prove a central limit theorem for it under regularity conditions. However, the sub-Gaussian parameter and the sub-Gaussian norm (corresponding to $\psi_2(x) = e^{|x|^2} - 1$) are \textit{not} equivalent: \cite{leskel2026sharp} proved that $0.612\approx \sqrt{3/8} \leq \xi_*/\|X\|_{\psi_2} \leq \sqrt{\log 2} \approx 0.832$ and showed the bounds to be sharp. Also, while $\xi_*^2$ does not generally have a closed form, \cite{atouani2025optimalsubgaussianvarianceproxy} studied the function $g_Y(\lambda; \sigma^2) = \frac12\lambda^2\sigma^2 - \log \E[e^{\lambda X}]$ to give explicit expressions of $\xi_*^2$ for three-point distributions. Other studies on the empirical moment generating function include \cite{Stewart2003-cd}, which focus on the supremum over $\lambda$ of the normalized process $n^{-1/2}\sum_{i=1}^n e^{\lambda X_i - \psi(2\lambda)/2} - \mathbb E[e^{\lambda X - \psi(2\lambda)/2}]$ where $\psi(\lambda) = \log \E[e^{\lambda X}]$.

\noindent \textbf{Notation}: If $I \subset \R$, then we set $I/2 := \{\lambda/2: \lambda \in I\}$. We denote $[n] = \{1,2,\dots,n\}$. We do not take $\N$ to include $0$, and will use the notation $\N_0 := \N \cup \{0\}$. 

\section{An Estimator for the sub-Gaussian Parameter}

For a sub-Gaussian distribution $P$ on $\mathbb{R}$, we let $\xi^2_*(P)$ denote its sub-Gaussian parameter. For $\Xi > 0$, we define the class of distributions
\[
 \mathcal{P}(\Xi) := \{ P \,:\, P \text{ sub-Gaussian}, \, \xi_*^2(P) \leq \Xi \}.
\]
To motivate our estimator, we begin with the following elementary observation about $\xi_*^2$. 
\begin{proposition}
\label{prop:xi-L-char}
    If $X \sim P$ is sub-Gaussian with variance $\sigma^2$, then 
    \[
    \xi_*^2(P) = \sup_{\lambda \in \R} L(\lambda ; P),\, \quad \text{ where } 
    L(\lambda; P) := \begin{cases}
        \frac{2}{\lambda^{2}} \log \E_P [e^{\lambda X}] & \lambda \ne 0 \\
        \sigma^2 & \lambda = 0
    \end{cases}
    \]
    is a continuous function. When $P$ is fixed, we will write $L(\lambda) \equiv L(\lambda; P)$.
\end{proposition}

Based on Proposition~\ref{prop:xi-L-char}, an immediate idea for estimating $\xi_*^2$ is to replace $L$ by its empirical analogue $L_n(\lambda) = \frac{2}{\lambda^2} \log \frac{1}{n}\sum_{i=1}^n e^{\lambda (X_i - \bar{X})}$ for $\lambda \neq 0$ and $L_n(0) = \frac{1}{n} \sum_{i=1}^n (X_i - \bar{X})^2$. Centering the empirical distribution by subtracting the sample mean $\bar{X}$ is critical even if $X_i$ is already mean-zero, otherwise $|L_n(\lambda)| \rightarrow \infty$ as $|\lambda| \rightarrow 0$. We may then consider estimating $\xi_*^2$ via $\sup_{\lambda \in \mathbb{R}} L_n(\lambda)$.  This however behaves poorly due to high variance of $L_n(\lambda)$ for large $\lambda$. In this work, we instead study the truncated maximization of $L_n$. We let $C_n > 0$ be a slowly diverging sequence and define our estimator
\[
\hat{\xi}_n^2 := \sup_{|\lambda| \leq C_n} L_n(\lambda). 
\]
We will see in our theoretical analysis that $C_n$ can be taken as $(\log n)^\alpha$ for any $\alpha \in (0, 1/2)$. 


\begin{figure}
    \centering
    \resizebox{0.6\linewidth}{!}{\input{figures/log-log_plot.pgf}}
    \caption{\textbf{Log-log plots for the mean absolute deviation for truncated $\hat \xi^2_n$ and untruncated $\hat \xi^2_{n, U}$.} }
   \label{fig:trunc-untrunc-comparison}
\end{figure}


To illustrate the importance of truncation, we conduct the following numerical experiment: we draw $X_1, \dots, X_n \sim N(0, 1)$ and compute the unconstrained estimator 
$\hat \xi_{n, U}^{2} := \sup_{\lambda \in \R} L_n(\lambda)$. We also compare it against $\hat \xi_n^2$ with $C_n = (\log n)^{1/4}$. In Figure~\ref{fig:trunc-untrunc-comparison}, we display the log-mean absolute deviation for $\hat \xi_n^2$ and $\hat \xi^2_{n, U}$ at different values of $n$ the averaged over 600 simulations; we take $n \in \{2^{j} \times 10^{2}: 2 \leq j \leq 7\}$. We observe that the error curve for $\hat \xi_n^2$ has a slope of approximately $-1/2$ in contrast to a slope of approximately $-1/4$ for $\hat \xi_{n, U}^2$. This is because $L_n(\lambda)$ exhibits unstable behavior when $\lambda$ is of order $\sqrt{\log n}$, slowing the convergence. Truncating removes this issue. 

\subsection{Convergence of the empirical moment generating function}

In order to study the consistency of $\hat \xi^2_n$ and quantify the rate of convergence, we must establish results on the uniform convergence of $L_n$ to $L$. To do this, we first study the empirical moment generating function. We shall denote the moment generating function and its empirical analogues by  
\[
M(\lambda; P) := \E_P[e^{\lambda X}], \quad M_n(\lambda) = \frac1n\sum_{i=1}^n e^{\lambda (X_i - \bar X)}, \quad \tilde M_n(\lambda) = \frac1n\sum_{i=1}^n e^{\lambda X_i}. 
\]
Similarly, we shall write $\psi(\lambda; P) := \log M(\lambda)$, $\psi_n(\lambda) = \log M_n(\lambda)$, and $\tilde \psi_n(\lambda) = \log \tilde M_n(\lambda)$. When $P$ is fixed, we will abbreviate $M(\lambda) \equiv M(\lambda ; P)$ and $\psi(\lambda) \equiv \psi(\lambda; P)$.


\begin{proposition}[\cite{feuerverger1989empirical}, Theorems 2.1 and 2.4]
\label{prop:feuerverger-uniform-convergence}
    Let $I$ be the largest interval on which $M(\lambda)$ is finite, and define $\tilde \psi_n(\lambda) = \log \tilde M_n(\lambda)$. Let $k \in \N$, and $K$ be a compact interval. Let $\tilde{\psi}^{(k)}$ be the $k$-th derivative of $\tilde{\psi}$ and likewise for $\tilde{\psi}_n^{(k)}$. 
    \begin{enumerate}
        \item[(a)] If $K\subset I$, then  $\sup_{\lambda \in K} |\tilde \psi_n^{(k)}(\lambda) - \tilde \psi^{(k)}(\lambda)| \xrightarrow{\mathrm{a.s.}} 0$ as $n \to \infty$. 
        \item[(b)] If $K \subset I/2$, then the empirical process $\sqrt n(\tilde \psi_n^{(k)}(\lambda) - \tilde \psi^{(k)}(\lambda))$ indexed by $\lambda$ converges weakly to a mean-zero Gaussian process on $K$. 
    \end{enumerate}
\end{proposition}
Since our truncation level $C_n$ is diverging (albeit at a slow rate), we need control over the convergence of $\tilde \psi_n$ and its derivatives beyond a fixed compact set. The next theorem establishes this for the empirical moment generating function.
\begin{theorem}
\label{prop:sub-gaussian-convergence-lemma}
    Let $k \in \N_0$.
    Fix $\alpha \in (0, 1/2)$, and set $C_n := (\log n)^\alpha$. Let $P$ be a sub-Gaussian distribution. Then for all $\eps > 0$,
    \[
    \sup_{\lambda \in [-C_n, C_n]} |M_n^{(k)}(\lambda) - M^{(k)}(\lambda; P)| = O_p(n^{-1/2 + \eps}),
    \]
    and this convergence is uniform over $P \in \mathcal P(\Xi)$. 
\end{theorem}
Having established a result on the convergence of the empirical moment generating function, we can easily extend the result to cumulant generating function. 
\begin{corollary}
\label{cor:cgf-unif-conv}
    Let $k \in \N_0$. Fix $\alpha \in (0, 1/2)$, and set $C_n := (\log n)^\alpha$. Let $P$ be a sub-Gaussian distribution. Then for all $\eps > 0$, 
    \[
    \sup_{\lambda \in [-C_n, C_n]}|\psi_n^{(k)}
(\lambda) - \psi^{(k)}(\lambda; P)| = O_p(n^{-1/2 + \eps}),   \]
and this convergence is uniform over $P \in \mathcal P(\Xi)$. 
\end{corollary}
Finally, control of the cumulant generating function extends directly to the $L_n$ function that we are maximizing. 
\begin{proposition}
\label{prop:L_n-cgf-bound}
    Let $C_n > 0$. For each $k \in \N_0$, $L_n$ and $L$ are differentiable everywhere on $\R$, and 
    \[
    \sup_{\lambda \in [-C_n, C_n]} |L_n^{(k)}(\lambda) - L^{(k)}(\lambda)| \leq \sup_{\lambda \in [-C_n, C_n]} |\psi^{(k+2)}_n(\lambda) - \psi^{(k+2)}(\lambda)|.
    \]
    Consequently, 
    \begin{enumerate}
        \item[(a)] If $C > 0$ is fixed, $\displaystyle \sup_{\lambda \in [-C, C]} |L_n^{(k)}(\lambda) - L^{(k)}(\lambda)| = O_p(n^{-1/2})$ uniformly over $P \in \mathcal{P}(\Xi)$. 
        \item[(b)] If $C_n = (\log n)^\alpha$ for $\alpha \in (0, 1/2)$, then for any $\eps > 0$, uniformly over $P \in \mathcal P(\Xi)$,
        \[
         \sup_{\lambda \in [-C_n, C_n]} |L_n^{(k)}(\lambda) - L^{(k)}(\lambda; P)| = O_p(n^{-1/2 + \eps}).
        \]
    \end{enumerate}
\end{proposition}


%Actually, the convergence in Proposition~\ref{prop:L_n-cgf-bound}(a) is also uniform over $\mathcal P(\Xi)$. The proof requires some machinery and is therefore delayed to Section~\ref{sec:minimax-opt}.
\subsection{Convergence of $\hat \xi_n^2$}
\label{sec:estimator-convergence}
Equipped with the above propositions, we are ready to prove the main result of this section. Let us define 
\[
\delta(C) \equiv \delta_P(C) := \sup_{|\lambda| \geq C} L(\lambda; P) - \sup_{|\lambda| \leq C} L(\lambda; P),
\]
and observe that $\delta(C)$ is a decreasing function with $\lim_{C \to \infty} \delta(C) \leq 0$.  

As we will show in Section~\ref{sec:minimax-opt}, the difficulty of this estimation problem is fundamentally related to the tail behavior of $\delta(C)$.
\begin{theorem}
\label{prop:second-consistency}
    Let $P$ be a sub-Gaussian distribution and let $X_1, \ldots, X_n$ be i.i.d samples from $P$. Then 
    \[
    |\hat \xi^2_n - \xi_*^2| \leq \sup_{|\lambda| \leq C_n} |L_n(\lambda) - L(\lambda)| + \delta_P(C_n) \vee 0.
    \]
    Consequently, if $C_n = (\log n)^\alpha$ for any $\alpha \in (0, 1/2)$, then 
    \begin{enumerate}
        \item[(a)] $\hat \xi_n^2 \xrightarrow{p} \xi_*^2$; \item[(b)] If there is $C_0 > 0$  such that $\delta(C_0) \leq 0$, then $\hat \xi_n^2 - \xi_*^2 = O_p(n^{-1/2 + \eps})$ for all $\eps > 0$;
        \item[(c)] If $\delta(C_0) < 0$ then $\hat \xi^2_n - \xi^2_* = O_p(n^{-1/2})$.
    \end{enumerate}
\end{theorem}
The proof of Theorem~\ref{prop:second-consistency} reveals that, so long as $\delta(C_0) = 0$ for some $C_0$, an exact rate of $n^{-1/2}$ can be achieved by choosing $C_n \equiv C_0$. For example, if $X \sim N(0, 1)$, then any such $C_0$ would work. However, we do not know $C_0$ a priori, so we cannot make this choice.
\begin{example}
    
We now give some cases when the guarantees of Theorem~\ref{prop:second-consistency} apply, or not. 
\begin{enumerate}
    \item[(a)] Suppose $X$ is bounded almost surely. Then $L$ obeys the inequality $L(\lambda) \leq 2\lambda^{-2}\log \E[e^{\|X\|_\infty \lambda}] = 2\|X\|_\infty \lambda^{-1}$,
        which vanishes as $\lambda \to \infty$. Therefore $\lim_{C \to \infty}\delta(C) = -\xi_*^2$, so Theorem~\ref{prop:second-consistency}(c) applies. If in particular $X \sim \text{Unif}[-1 ,1]$, then $L$ satisfies the hypotheses of Proposition~\ref{prop:ptwise-asymptotic-normality} so $\sqrt n(\hat \xi_n^2 - \xi_*^2) \xrightarrow{d}  N(0, 4/45)$.
        
\item[(b)] If $X \sim N(0, 1)$, then $L(\lambda) \equiv 1$ for all $\lambda \in \R$, so $\delta(C) \equiv 0$. Therefore (c) does not apply, but (b) does.

\item[(c)] If $X \sim \frac12N(0, 1) + \frac12N(0, 2)$, then $\delta(C) > 0$ for all $C > 0$ so only Theorem~\ref{prop:second-consistency}(a) holds.
\end{enumerate}
\end{example}
\begin{figure}[h]
    \centering
    \resizebox{0.45\linewidth}{!}{\input{figures/L_plots.pgf}}
    \caption{\textbf{Plots of $L(\lambda; P)$.} In Example 1(a), $\argmax L$ exists so $\delta$ is eventually negative. In (c), $L$ is increasing away from 0 so that $\delta(C) > 0$ for all $C > 0$.}
    \label{fig:placeholder}
\end{figure}

\subsection{Inference}
\label{sec:inference}
Under even stronger assumptions, we can also derive the limiting distribution of $\hat \xi_n^2$. 
\begin{proposition}
\label{prop:ptwise-asymptotic-normality}
For $\lambda \ne 0$, define   
    \[
    V(\lambda) := \frac{4}{\lambda^4}\left(\frac{M(2\lambda)}{M(\lambda)^2} - 2\lambda\psi'(\lambda)+ \lambda^2\sigma^2 - 1\right).
    \]
    Also, let $V(0) := \E X^4 - \sigma^4$. Suppose that $L$ is uniquely maximized at $\lambda^*$, $L''(\lambda^*) < 0$, $V(\lambda^*) > 0$, and there is $C_0$ such that $\delta(C_0) < 0$. 
    Then, for any $C_n \to \infty$, we have 
    \[
    \sqrt n(\hat \xi_n^2 - \xi_*^2) \xrightarrow{d} N(0, V(\lambda^*)).
    \]
\end{proposition}
The hypothesis that $L$ be uniquely maximized is necessary: if $X \sim N(0, 1)$, then $L(\lambda) \equiv 1$ so that all of $\R$ maximizes $L$. In this case, the maximizers of $L_n$ will be entirely determined by stochastic fluctuations in $L_n$, so the limiting distribution is not Gaussian. We include the $N(0, 1)$ case in Figure~\ref{fig:asymptotic-normality} to explicitly illustrate this behavior. 

According to Proposition~\ref{prop:ptwise-asymptotic-normality}, we may construct asymptotically valid confidence intervals by plugging in empirical estimates of the quantities involved. Let 
\[
\hat V_n(\lambda) := \frac{4}{\lambda^4}\left(\frac{M_n(2\lambda)}{M_n(\lambda)^2} - 2\lambda\psi_n'(\lambda) + \lambda^2\hat\sigma^2_n - 1\right),
\]
where $\hat \sigma_n^2$ is any consistent estimator of $\sigma^2$. Then, if $z_{\alpha/2}$ denotes the $(1 - \alpha/2)$-quantile of the standard normal distribution and $\lambda^*_n$ denotes a maximizer of $L_n$ over $[-C_n, C_n]$, it follows that 
\[
\left[\hat \xi_n^2 -  z_{\alpha/2}\sqrt{\frac{\hat V_n(\lambda^*_n)}{n}}, \hat \xi^2_n + z_{\alpha/2}\sqrt{\frac{\hat V_n(\lambda^*_n)}{n}}\, \right]
\]
is an asymptotic $(1 - \alpha)$-level confidence interval for $\xi_*^2$. 

To verify Proposition~\ref{prop:ptwise-asymptotic-normality}, we generate samples $X_1, \dots, X_n \sim P$ for the following choices of $P$: (a) $\mathrm{Unif} [-1, 1]$,
(b) $\frac{1}{3}\delta_{-2} + \frac{2}{3}\delta_1$, 
(c) $\frac12 \mathrm{Unif}[-3, 3] + \frac12N(0, 1)$, 
(d) $N(0, 1)$.

We set $C_n = (\log n)^{1/4}$, $n = 1000$, and run 600 simulations. For each simulation, we compute the estimated asymptotic variance $\hat V_n(\lambda^*_n)$ and take the average over all simultations and compare it to the sample variance.

\begin{figure}[h]
\centering
    \resizebox{0.6\linewidth}{!}{\input{figures/asymptotic_normality.pgf}}
\caption{\textbf{Distribution of $\sqrt n(\hat \xi_n^2 - \xi_*^2)$ for four choices of $P$.} Proposition~\ref{prop:ptwise-asymptotic-normality} gives asymptotic normality for all cases but $N(0, 1)$.}
    \label{fig:asymptotic-normality}
\end{figure}
In practice, the hypotheses of Theorem~\ref{prop:second-consistency} and Proposition~\ref{prop:ptwise-asymptotic-normality} are difficult to check. One heuristic approach is to plot $L_n$ on $[-C_n, C_n]$, and check if the maximizers are in the interior. An affirmative answer would be consistent with the hypotheses in Theorem~\ref{prop:second-consistency}(c) and Proposition~\ref{prop:ptwise-asymptotic-normality}; on the other hand if $L_n$ is flat or monotone then the guarantees likely do not apply and the confidence interval may not be reliable.

\section{Minimax Risk of Estimating $\xi_*^2$}

In the previous section, we saw that the tail behavior of the truncation gap $\delta_P$ determines the rate at which $\hat \xi^2_n$ converges. In Section~\ref{sec:minimax-lower}, we show that $\delta_P$ governs the difficulty of estimating $\xi_*^2$ itself, regardless of the estimator involved. To do this, we consider three regimes: first, when no uniform assumptions are made on $\delta_P$, we show that the minimax risk is $\Omega(1)$. If a uniform polynomial decay is imposed, then the risk becomes $\Omega(1/\log n)$, and if the $\delta_P$'s eventually become zero or negative, the risk becomes $\Omega(n^{-1/2})$. For the third regime, we show that $\hat \xi_n^2$ matches the $n^{-1/2}$-rate, making it minimax optimal.

\subsection{Lower bounds on the minimax risk}
\label{sec:minimax-lower}
This subsection characterizes the difficulty of the estimation problem in terms of an $r$ function that determines the uniform rate of decay for the $\delta_P$'s. Intuitively, slower rates of decay on $\delta_P$ correspond to harder estimation problems because the supremum of $L$ is more difficult to approximate, so any estimator must search for maximizers for an increasingly large range of $\lambda$.

To make this precise, let us write
\[
\mathcal P_c(\Xi) := \{P \in \mathcal P(\Xi): \text{$P$ is compactly supported}\},
\]
and given any $C_0 > 0$ and decreasing function $r: (0, \infty) \to \R$, define the classes
\[
\mathcal P(\Xi, C_0, r) := \{P \in \mathcal P(\Xi): \delta_P(C) \leq r(C) \text{ for all $C \geq C_0$}\}
\]
and 
\[
\mathcal P_c(\Xi, C_0, r) := \{P \in \mathcal P_c(\Xi): \delta_P(C) \leq r(C) \text{ for all $C \geq C_0$}\}.
\]

The following proposition characterizes membership in $\mathcal P_c(\Xi, C_0, r)$.
\begin{proposition}
\label{prop:unif-supp-delta-inclusion}
    Let $P \in \mathcal P(\Xi)$. If $P$ is supported on $[-1, 1]$, then $P \in \mathcal P(\Xi, 0, r_0)$ where $r_{0}(t) = 2/t$. If in addition $\sigma^2(P) \geq \frac32\delta_0$, then $P \in \mathcal P(\Xi, 8\delta_0^{-1}, r_-)$ where $r_-(t) = - \delta_0$.
\end{proposition}
Let us now define, for any $\gamma \in [0, 1/2]$, 
\[
 r_\gamma(t) = 2^{\frac{1 - 2\gamma}{1 - \gamma}} t^{\frac{2\gamma - 1}{1 - \gamma}}, \quad t > 0.
\]
There are two values of $\gamma$ that are particularly important. When $\gamma = 1/2$, $r_\gamma(t) \equiv 1$, so that no rate of decay is imposed on $\delta_P$. Meanwhile, for $\gamma = 0$, we have $r_\gamma(t) = 2/t$, the rate of decay that corresponds to random variables supported on $[-1, 1]$, as in Proposition~\ref{prop:unif-supp-delta-inclusion}.

Let $0 \leq \gamma_1 \leq \gamma_2 \leq 1/2$. It then follows that $r_{\gamma_1}(t) \leq r_{\gamma_2}(t)$ for any $t \geq 2$, and therefore we have the inclusion 
\[
\mathcal P(\Xi, 2, r_{\gamma_1}) \subset \mathcal P(\Xi, 2, r_{\gamma_2}).
\]
Thus as $\gamma$ increases, $\mathcal P(\Xi, 2, r_\gamma)$ grows larger and estimation becomes harder. The next proposition quantifies this.



\begin{proposition}
\label{prop:minimax-rates}
There exists a universal constant $c_0> 0$ such that: 
    \begin{enumerate}
        \item[(a)] For any $\gamma \in [0, 1/2]$, we have 
        \[
        \inf_{\hat \xi^2} \sup_{P \in \mathcal P_c(\Xi, 0, r_\gamma)} \E_P|\hat \xi^2(X_1, \dots, X_n) - \xi_*^2(P)| \geq \frac{c_0 \cdot \Xi }{(\log 2n)^{1 - 2\gamma}}
        \]
        and in particular, 
        \[
        \inf_{\hat \xi^2} \sup_{P \in \mathcal P_c(\Xi)} \E_P|\hat \xi^2(X_1, \dots, X_n) - \xi_*^2(P)| \geq c_0 \cdot \Xi.
        \]
        \item[(b)] Let $\delta_0 \in [0, 1/2]$ and set $r_{-}(t) = -\delta_0$; let $C_0 \geq 16$. Then 
        \[
        \inf_{\hat \xi^2} \sup_{P \in \mathcal P_c(\Xi,C_0, r_-)} \E_P|\hat \xi^2(X_1, \dots, X_n) - \xi_*^2(P)| \geq \frac{c_0 \cdot \Xi}{\sqrt n}.
        \]
    \end{enumerate}
\end{proposition}


\subsection{Upper bounds on the minimax risk via $\hat \xi_n^2$}
\label{sec:minimax-opt}
In Theorem~\ref{prop:second-consistency}(a), we proved the pointwise consistency of $\hat \xi_n^2$. If the tail of $\delta_P$ can be uniformly controlled, then we may also get uniform consistency.
\begin{proposition}
\label{prop:P(Xi,K)-unif-consistency}
    Set $C_n = (\log n)^{\alpha}$ for any $\alpha \in (0, 1/2)$. For any $C_0 > 0$ and decreasing function $r: (0, \infty) \to \R$ with $\lim_{t \to \infty} r(t) = 0$, we have $|\hat \xi_n^2 - \xi_*^2(P)|\leq  r(C_n) + O_p(n^{-1/2 + \eps}) = o_p(1)$
    uniformly over $P \in \mathcal P(\Xi,C_0, r)$.
\end{proposition}

By Proposition~\ref{prop:P(Xi,K)-unif-consistency}, we have that $\sup_{P \in \mathcal{P}(\Xi, 0, r_{\gamma})} \mathbb{E}_P |\hat{\xi}^2_n - \xi^2_*(P)| \leq O( (\log n)^{\frac{\alpha(2\gamma - 1)}{1 - \gamma}})$ when $C_n$ is chosen as $(\log n)^\alpha$ for $\alpha \in (0, 1/2)$ and $\gamma \in [0, 1/2]$. This does not match the lower bound of $(\log 2n)^{2 \gamma - 1}$ even in the limit $\alpha \rightarrow 1/2$. It is an open question whether to improve the lower bound or the upper bound in order to close the gap.

The remainder of this section works towards a proof that $\hat \xi_n^2$ achieves the rate in Proposition~\ref{prop:minimax-rates}(c), uniformly over $\mathcal P(\Xi, C_0, r_-)$. First, we show that the maximizers of $L_n$ and $L$ lie inside a compact interval $[-C_0, C_0]$ with probability tending to 1, uniformly over $\mathcal P(\Xi, C_0, r_-)$.
\begin{proposition}
\label{prop:maximizer-uniform-bound}
    Set $C_n = (\log n)^\alpha$ for any $\alpha \in (0, 1/2)$. Let $\Lambda^*_n$ denote the maximizers of $L_n$ in $[-C_n, C_n]$ and let $\Lambda^*(P)$ denote the maximizers of $L$. Then   
\[
    \limsup_{n\to\infty} \sup_{P \in \mathcal P(\Xi, C_0, r_-)} \P\left(\Lambda^*_n \cup \Lambda^*(P) \not\subset [-C_0, C_0]\right) = 0.
    \]
\end{proposition}
Proposition~\ref{prop:maximizer-uniform-bound} says that, with high probability, we can identify $L_n$ with its restriction to $[-C_0, C_0]$, independent of $P \in \mathcal P(\Xi, C_0, r_-)$, and similarly for $L$. 


We shall use this to gain control over empirical moment generating function. To that end, define
\[
\mathcal F_k = \{x \mapsto x^ke^{\lambda x}: \lambda \in [-C_0, C_0]\} \subset \R^\R.
\]
Informally, the next proposition gives an empirical process central limit theorem for $\mathcal F_k$, and guarantees that the convergence is uniform over $\mathcal P(\Xi)$. Precise definitions are given in Section~\ref{sec:uniform-donsker}. 
\begin{proposition}
\label{prop:mgf-uniform-donsker}
    $\mathcal F_k$ is Donsker and pre-Gaussian, both uniformly over $\mathcal P(\Xi)$. 
\end{proposition}
Using Proposition~\ref{prop:mgf-uniform-donsker}, we can prove the following, which is a re-statement of Proposition~\ref{prop:L_n-cgf-bound}(a). 
\begin{proposition}
\label{prop:L_n-uniform-rootn}
% {\color{orange}
% Fix $C_0 > 0$ and $k \in \N_0$. The empirical process $\lambda \mapsto \sqrt{n} \{L_n^{(k)}(\lambda) - L^{(k)}(\lambda; P)\}$ converges to a mean-zero Gaussian process on $[-C_0, C_0]$, uniformly for all $P \in \mathcal{P}(\Xi)$. As a direct consequence, as $n \to \infty$, uniformly over $P \in \mathcal P(\Xi)$.
%     \[
%      \sup_{|\lambda| \leq C_0} |L_n^{(k)}(\lambda) - L^{(k)}(\lambda; P)| = O_p(n^{-1/2}),
%     \]
% }
Fix $k \in \N_0$. As $n \to \infty$, uniformly over $P \in \mathcal P(\Xi)$,
    \[
     \sup_{|\lambda| \leq C_0} |L_n^{(k)}(\lambda) - L^{(k)}(\lambda; P)| = O_p(n^{-1/2}).
    \]
\end{proposition}


We are now ready to make the convergence in Theorem~\ref{prop:second-consistency}(c) uniform over $\mathcal P(\Xi, C_0, r_-)$. 
\begin{theorem}
\label{prop:hat xi-minimax-optimality}
    Let $C_n \to \infty$. As $n \to \infty$, uniformly over $\mathcal P(\Xi, C_0, r_-)$, $\hat \xi_n^2 - \xi_*^2(P) = O_p(n^{-1/2})$.
\end{theorem}
Combined with Proposition~\ref{prop:minimax-rates}(b), this shows that $\hat \xi^2_n$ achieves the optimal rate. Because $C_n$ can be any sequence that diverges, we have also shown that $\hat \xi^2_n$ is adaptive in the sense that it does not require knowledge of $\Xi$, $C_0$, or $\delta_0$.
\section{Rate of Divergence Under Violation of sub-Gaussianity}
We now study the mis-specified case when the underlying distribution is not sub-Gaussian. In this setting, it is desirable for $\hat \xi^2_n$ to exhibit clear signs of divergence, which 
would serve as qualitative evidence against the sub-Gaussian assumption. As we next 
show, $\hat\xi^2_n$ diverges almost surely whenever sub-Gaussianity fails.

\begin{proposition}
\label{prop:divergence}
Let $X$ be a non-sub-Gaussian random variable, i.e. if $\mathbb{E} e^{\lambda X} = \infty$ for some $\lambda \in \mathbb{R}$ or if for any $\xi^2 > 0$, there exists $\lambda \in \mathbb{R}$ such that $\mathbb{E} e^{\lambda X} \geq e^{ \lambda^2 \xi^2 / 2}$. Let $X_1, \ldots, X_n$ be i.i.d. copies of $X$. If $C_n \rightarrow \infty$, then $\hat{\xi}^2_n \xrightarrow{\mathrm{a.s.}} \infty$. 
\end{proposition}

Although Proposition~\ref{prop:divergence} guarantees that $\hat{\xi}_n^2$ will go to infinity, we still expect $\hat{\xi}^2_n$ to diverge faster when $X$ has a heavier tail and thus a stronger departure from sub-Gaussianity. To quantify the rate of divergence, we state a Proposition which gives a deterministic lower bound on $\hat{\xi}_n^2$ in terms of the range of the data. 

\begin{proposition}
\label{prop:divergence2}
Let $X_1, \ldots, X_n$ be any collection of real numbers and let $\Delta_n := \max_{i = 1, \ldots, n} | X_i - \bar{X}|$. If $\Delta_n \geq \frac{2\log n}{C_n}$, then $\hat{\xi}^2_n \geq \frac{\Delta_n^2}{2 \log n}$. 
\end{proposition}

If we set $C_n = (\log n)^\alpha$ for an $\alpha \in (0, 1/2)$, then Proposition~\ref{prop:divergence2} applies when $\Delta_n \geq (\log n)^{1- \alpha}$. Since $\Delta_n = O_p(\sqrt{ \log n})$ for sub-Gaussian data (\cite{Wainwright_2019}, Exercise 2.12), we expect $\Delta_n$ to be of larger order when $X_1, \ldots, X_n$ are not sub-Gaussian. We consider some examples:
\begin{enumerate}[leftmargin=3mm]
\item Suppose $X_1, \ldots, X_n$ are distributed according to the mean-zero Laplace distribution. Then it holds that $\Delta_n$ is of order $\Theta_p(\log n)$ so that the condition of Proposition~\ref{prop:divergence2} holds for all large enough $n$ and we have that $\hat{\xi}^2_n$ diverges at rate $\Omega_p(\log n)$. 
\item Suppose $X_1, \ldots, X_n$ have the heavy tail t-distribution with degree $\nu > 0$. Recall that the expectation exists only when $\nu > 1$ and the variance is finite only when $\nu > 2$; the t-distribution coincides with the Cauchy distribution when $\nu = 1$. In this case, we have that $|X_{(n)} - X_{(1)}|$ is of order $\Theta_p( n^{1/\nu} )$ so that $\hat{\xi}^2_n$ diverges at rate $\Omega_p( \frac{n^{2/\nu} }{\log n})$. 
\end{enumerate}

%Part (b) quantifies the rate of divergence directly in terms of the observed data, and the hypothesis is easy to verify. If $C_n = (\log n)^\alpha$ for $\alpha \in (0, 1/2)$, it requires that $X_{(n)} - \bar X = \Omega_p((\log n)^{1 - \alpha})$, which encompasses any non sub-Gaussian distribution as $\alpha \to 1/2$ (recall that the sample maximum of a sub-Gaussian distribution grows at rate $(\log n)^{1/2}$).

%For example, if $X$ is exponentially distributed, then $X_{(n)} = \Theta_p(\log n)$, so $\hat \xi^2_n = \Omega_p(\log n)$. For heavier-tailed distributions where $X_{(n)} = \Omega_p(n^c)$, we have $\hat \xi^2_n = \Omega_p(n^{2c}(\log n)^{-1})$; for instance if $X$ is standard Cauchy, then $X_{(n)} = \Omega_p(n)$ and $\hat \xi^2_n = \Omega_p(n^{2}(\log n)^{-1})$.
\section{Application to Gene Ontology Enrichment Studies}
\label{sec:go-enrichment}
In this section, we apply the sub-Gaussian estimator to the setting of large scale permutation tests in Gene Ontology (GO) enrichment analysis, a multiple testing problem with large numbers of hypotheses relative to the number of permutations. In this case, the empirical p-value is not granular enough; we propose using the sub-Gaussian tail bound instead.

\noindent \textbf{Background}:
A standard goal in disease gene identification is to determine whether candidate genes obtained from high-throughput sequencing experiments are enriched for certain GO terms, which represent biological processes.

One method to do this \citep{levi2021domino, liu2025beyond} is as follows: gene scores representing significance are placed onto a gene-gene interaction (GGI) network. Then, regions exhibiting both high connectivity and candidate gene enrichment called \emph{modules} are identified using \emph{active module identification (AMI) algorithms}. Each (module, GO term) pair is then tested via the hypergeometric test, and the resulting negative log-10 p-value is called the \emph{HG-enrichment score}. The \emph{enrichment} $Z^{\text{obs}}_g$ is the maximal enrichment across all modules, and if $Z^{\text{obs}}_g$ is greater than some threshold, $g$ is called \emph{HG-significant}.

To test the statistical significance of the $Z^{\text{obs}}_g$'s, a standard approach are permutation tests: generate $n$ permutations of the original data, rerun the full pipeline on each, and obtain scores $Z^{(i)}_g$ for each $g \in [G]$ and $i \in [n]$. HG-significant GO terms whose empirical p-value passes Benjamini-Hochberg (BH) correction at level $\alpha$ are called \emph{empirically validated.}

The empirical p-value obtained by this test is no smaller than $n^{-1}$. However, $n$ is typically much smaller than $G$ because obtaining permutation requires rerunning the full module detection pipeline. As a result, after multiple testing correction such as BH, it can be impossible to reject any GO terms. This motivates new procedures to generate valid p-values for each GO term $g$ that are fine-grained enough to pass false discovery rate corrections. 

\noindent \textbf{The sub-Gaussian transform}: 
Fix a GO term $g \in [G]$. If $Z^{(i)}_g$ is a sub-Gaussian distribution with parameter $\xi_*^2(g)$, then the sub-Gaussian tail bound would justify the use of the p-value 
\[
\mathbf{p}_g^{\text{sG}} := \exp\biggl(-\frac{(Z^{\text{obs}}_g - \bar Z_g)^2}{2\hat \xi^2_{n}(g)}\biggr)
\]
where $\bar Z_g = \frac1n\sum_{i=1}^n Z_g^{(i)}$ and $\hat \xi^2_{n}(g)$ is the value of the estimator given $Z^{(1)}_g,\dots, Z^{(n)}_g$. 
Such an assumption can be made in this context because GO enrichment scores are bounded (the largest enrichment score attainable is $\log_{10}{N \choose N/2}$ where $N$ is the number of background genes), and therefore sub-Gaussian. However, the sub-Gaussian tail inequality using this bound on the distribution is too loose, and it is necessary to use a more precise estimator like $\hat \xi^2_n(g)$.

\noindent \textbf{The Peaks-over-thresholds method}: 
Another approach \citep{@knijnenburg2009fewer} is to define the \textit{excess function}

\[
F_g^{(u)}(y) =  \frac{ F_g(u+y) - F_g(u)}{1 - F_g(u)} = \mathbb P(Z_g - u \le y  \, | \, Z_g > u)
\]
where $F_g(u)$ is the distribution function of $Z_g$ under the null and $u \in \R$ is a threshold parameter. According to the Pickands-Balkema-de Haan theorem, as $u \to \sup \{x: F_g(x) < 1\}$, $F_g^{(u)}$ converges in distribution to a Generalized Pareto Distribution (GPD) with parameters $m_g, s_g, k_g$. This motivates fitting a GPD to the 
exceedances $\{Z^{(b)}_g - u_g : Z^{(b)}_g > u_g\}$, 
yielding estimates $\hat m_g$, $\hat s_g$, and $\hat k_g$, and 
the p-value

\[
\mathbf{p}_g^{\mathrm{POT}} := (1 - \hat F_g(u_g))
\left(1 +  
\frac{\hat k_g(Z^{\text{obs}}_g - u_g - \hat m_g)}{\hat s_g}
\right)^{-1/\hat k_g},
\]


where $\hat F_g$ is the empirical distribution function. In practice, $u_g$ must be chosen with care. If it is too large, there may be too few exceedances to obtain a reliable fit. A common choice for $u_g$ is the 90-quantile.

\noindent \textbf{Method Comparison}:
We next compare the GO terms rejected by both methods on widely used complex gene datasets. Following \cite{liu2025beyond}, we used the ``TNFa''  dataset \citep{schmidt2015acute} with the ``DIP" network  \citep{xenarios2000dip} and the ``Fly Transcriptome'' \citep{sun2024whole} dataset with the ``STRING" network \citep{szklarczyk2017string}. We  used the PAPER \citep{crane2024root}, DOMINO \citep{levi2021domino}, and FDRnet \citep{yang2021efficient} AMI algorithms. We excluded the ``Aneuploidy1'' and ``Aneuploidy2'' datasets because they exhibited very low numbers of HG-enriched GO terms, and similarly for the HotNet2 \citep{leiserson2015pan} algorithm.


For each algorithm and dataset, we compute $\mathbf{p}^{\mathrm{sG}}_g$ and $\mathbf{p}^{\mathrm{POT}}_g$ per GO term, apply BH correction at level $\alpha = 0.05$, and in Table~\ref{tab:ev-go-term} we report the number of empirically-validated GO terms followed by the number of BH-significant GO terms in parentheses. 


\begin{table}[h]
    \centering
    \begin{tabular}{ccc}
    \hline
        & sub-Gaussian & POT \\\hline
       PAPER & 181 (279) & 115 (138) \\\hline
       DOMINO & 137 (143) & 167 (169) \\\hline
       FDRnet & 1 (1) & 1 (3) \\\hline
    \end{tabular} $\quad$
    \begin{tabular}{ccc}
    \hline
        & sub-Gaussian & POT \\\hline
       PAPER & 0 (0) & 4 (10) \\\hline
       DOMINO & 5 (5) & 25 (58) \\\hline
       FDRnet & 13 (39) & 20 (82) \\\hline
    \end{tabular}
\caption{\textbf{Empirically validated GO terms} by AMI algorithm 
(rows) and p-value method (columns); BH-significant counts in 
parentheses. Left: Fly Transcriptome. Right: TNFa.}
    \label{tab:ev-go-term}
\end{table}

On the Fly Transcriptome dataset, the sub-Gaussian transform yields more empirically validated GO terms under PAPER and comparable results for DOMINO and FDRnet. On the TNFa dataset, POT achieves more rejections across all three algorithms. The relative performance therefore depends heavily on the null distribution, and the two methods are best understood as complementary.

The POT method is subject to different modes of failure. First, the rate of convergence can be slow and very distribution dependent. For example, if the underlying distribution is normal, the convergence of the empirical excess function $F^{(u)}_g$ to a GPD is of order $1/u^2$ as $u \to \infty$ \citep{Degen_Embrechts_Lambrigger_2007}, but the threshold $u_g$ can grow at most at the rate of the sample maximum, which is $O(\sqrt{\log n})$ for sub-Gaussian data. The rate is therefore $O(1/\log n)$, which is slow enough that in our case of permutation testing where $n$ is relatively small, it is not clear whether the asymptotic regime has been reached. In this case, the p-values produced may not be valid. Second, for PAPER on the Fly Transcriptome dataset, 255 out of 279 (91.4\%) of the POT fits have $\hat k_g < 0$, which corresponds to a distribution supported on $[0, m_g - s_g/k_g]$. Among these fits, 8 satisfied $Z_g^{\text{obs}} > m_g - s_g/k_g$ resulting in a p-value of 0, which is a clear underestimate. Third, the standard choice of $u_g$ as the 90th quantile does not diverge as $n \to \infty$, but taking it too grow too fast leads very few exceedances, particularly for light-tailed distributions in the sub-Gaussian case. Moreover, the problem of selecting this threshold $u_g$ does not have a general solution \citep{Degen_Embrechts_Lambrigger_2007}.


On the other hand, $\hat \xi_n^2$ may also exhibit slow convergence under certain distributions (for example $\frac12 N(0, 1) + \frac12 N(0, 2)$). If the underlying distribution is heavy tailed, the divergence of $\hat \xi_n^2$ would make the p-values useless in comparison to the POT approach. In practice, we view these methods as complementary; when the data is sufficiently light-tailed the sub-Gaussian transform may be more reliable, in other cases POT may be better.

\section{Discussion}

\label{sec:discussion}
We have studied the estimation of the sub-Gaussian parameter $\xi_*^2$ and proposed an estimator $\hat \xi_n^2$ based on the truncated maximization of the empirical analogue $L_n$. One of our main findings is that the truncation gap $\delta_P$ controls the fundamental difficulty in estimating $\xi_*^2$: without control over $\delta_P$ no estimator can be uniformly consistent; stronger decay assumptions on $\delta_P$ lead to faster minimax rates, culminating at a rate of $n^{-1/2}$ at which $\hat \xi_n^2$ is minimax optimal and adaptive. Under departures from sub-Gaussianity, the estimator diverges almost surely at a rate depending on the range of the observed data, providing a natural diagnostic of mis-specification. Our GO enrichment application demonstrate that for large scale permutation testing, sub-Gaussian tail bounds can act as a complementary approach to the POT method.


There are a few limitations to our work. For one, the hypotheses of Theorem~\ref{prop:second-consistency} and Proposition~\ref{prop:ptwise-asymptotic-normality} are difficult to deduce from observed data, though in practice the heuristic suggested in Section~\ref{sec:inference} may give a loose estimate of the regime. The rate in Theorem~\ref{prop:second-consistency}(b) has an additional $\eps$ term which is not known to be sharp, and a more detailed analysis is required to close this gap. While taking $C_n = (\log n)^\alpha$ is theoretically sound, a data-driven choice would be more practical. Finally, the minimax lower bound in Proposition~\ref{prop:minimax-rates}(b) is not met, and it is an open problem to determine the minimax rate when other conditions on $\delta$ are imposed. Addressing these gaps are natural directions for future work.




\bibliographystyle{dcu}
\bibliography{all}

\newpage

\appendix

\section{Appendix}
\label{sec:appendix}
\subsection{Proofs for Section 2}
\begin{proof}[Proof of Proposition~\ref{prop:xi-L-char}]
    Firstly, since $X$ is sub-Gaussian, $M(\lambda) < \infty$ for all $\lambda \in \mathbb{R}$ and hence $|L(\lambda)| < \infty$ as well. Now we observe that, by rearranging the inequality, we have that $\xi > 0$ satisfies
    \[
    \E[e^{\lambda X}] \leq \exp\left(\frac{\lambda^2\xi^2}{2}\right)
    \]
    if and only if $\xi^2 \geq L(\lambda)$ for $\lambda \ne 0$. Since $\xi^2 \geq \sigma^2$ and $\lambda$ was arbitrary, we have $\xi^2 \geq \sup_{\lambda \in \R} L(\lambda)$, so  $\xi^2_* \geq \sup_{\lambda \in \R} L(\lambda)$ too. Conversely, for any $\eps > 0$ there must be $\lambda \in \R$ such that 
    \[
    \E[e^{\lambda X}] >\exp\left(\frac{\lambda^2(\xi^* - \eps)^2}{2}\right),
    \]
    i.e. $(\xi^* - \eps)^2 < 2\lambda^{-2}\psi(\lambda) = L(\lambda)$, so that $\xi^*$ is the least upper bound. It remains to show that $L$ is continuous. This can be done by applying L'H\^opital's rule to $ \lim_{\lambda \to 0} 2\lambda^{-2}\psi(\lambda)$ and recalling that $\psi''(0) = \sigma^2$. 
\end{proof}
\subsubsection{Proof of Theorem~\ref{prop:sub-gaussian-convergence-lemma}}
Our first main goal of this section is prove Theorem~\ref{prop:sub-gaussian-convergence-lemma}. 

First, we need the following concentration inequality. Note that we cannot immediately apply any concentration inequalities to $M_n^{(k)}$, because the terms in the summation are not independent. 
\begin{proposition}
\label{propn:mgf-concentration-ineq}
    Let $k \in \N_0$. Assume that $M(\lambda)$ exists at $\lambda$, and that $|X| \leq B$. Then, for any $|\lambda| \leq C$, 
    \[
    \P\left(|M_n^{(k)}(\lambda) - M^{(k)}(\lambda)| > t\right) \leq 2\exp\left(\frac{-2n^{1/2}}{4B^2}\right) + 2\exp\left(\frac{-nt^2}{8B^ke^{4BC+1}}\right).
    \]
\end{proposition}
\begin{proof}
Assume that $B > 1$ without loss of generality. 
     Put $\tilde M_n(\lambda) := \frac1n\sum_{i=1}^n e^{\lambda X_i}$, and then observe that by the triangle inequality,
    \begin{align*}
        |M_n^{(k)}(\lambda) - M^{(k)}(\lambda)| &= |M_n^{(k)} - \tilde M_n^{(k)}(\lambda)| + |\tilde M_n^{(k)}(\lambda) - M^{(k)}(\lambda)|.
    \end{align*}
    Now by Hoeffding's inequality, we have, for any $\delta_n \to 0$, 
    \begin{equation}
        \P(|\bar X| > \delta_n) \leq 2\exp\left(-\frac{2n\delta_n^2}{(2B)^2}\right). \label{eqn:hoeffding-mean}
    \end{equation}
    Suppose that $|\bar X| < \delta_n$. For the first term, define $g(\mu) := \frac1n\sum_{i=1}^n (X_i - \mu)^k e^{\lambda (X_i - \mu)}$. Then $g(0) = \tilde M_n^{(k)}(\lambda)$ and $g(\bar X) = M_n^{(k)}(\lambda)$, so by the mean-value theorem there is $\tilde \mu \in [0 \wedge \bar X, 0 \vee \bar X]$ such that 
    \[
    M_n^{(k)}(\lambda) - \tilde M_n^{(k)}(\lambda) = g'(\tilde \mu)\bar X,
    \]
    but for any $\mu \in [-C, C]$, 
    \[
    |g'(\mu)| \leq \frac1n\sum_{i=1}^n \left|k(X_i - \mu)^{k-1} + \lambda(X_i - \mu)^{k}\right|e^{\lambda(X_i - \mu)} \leq \left(k(2B)^{k-1} + C(2B)^{k}\right)e^{2B|\lambda|},
    \]
    so 
    \[
    |M_n^{(k)}(\lambda) - \tilde M_n^{(k)}(\lambda)| \leq \left(k(2B)^{k-1} + C(2B)^{k+1}\right)e^{2B|\lambda| }\delta_n.
    \]
    Therefore sending $\delta_n \to 0$ makes this term smaller than $t/2$. 
    For the second term, note that when $\lambda \in [-C, C]$, 
    \[
    \tilde M^{(k)}_n(\lambda) = \frac1n\sum_{i=1}^n X_i^{k}e^{\lambda X_i} \leq B^ke^{B|\lambda|} \leq B^ke^{2BC},
    \]
    so by Hoeffding's inequality again 
    \begin{align}
        \P\left(|\tilde M_n^{(k)}(\lambda) - M^{(k)}(\lambda)| > \frac{t}{2e^{|\lambda|\delta_n}}\right) &\leq 2\exp\left(-\frac{2n}{(2B^ke^{2BC})^2}\frac{t^2}{4e^{|\lambda| \delta_n}}\right)\notag\\ &\leq  2\exp\left(-\frac{nt^2}{8B^{2k}e^{4BC+2|\lambda|\delta_n}}\right)
        \notag \\ 
        &\leq 
2\exp\left(-\frac{nt^2}{8B^{2k}e^{4BC+1}}\right)\label{eqn:hoeffding-mgf}
    \end{align}
    where the last inequality is valid for large enough $n$, since $\delta_n \to 0$. Setting $\delta_n := n^{-1/4}$ in (\ref{eqn:hoeffding-mean}) and (\ref{eqn:hoeffding-mgf}) yields the desired conclusion.
\end{proof}
Equipped with Proposition~\ref{propn:mgf-concentration-ineq}, we are ready to prove a weaker analogue of Theorem~\ref{prop:sub-gaussian-convergence-lemma}, first for bounded random variables. 
\begin{proposition}
\label{prop:bounded-convergence-lemma}
    Let $k \in \N_0$. Fix $\alpha, \beta \in (0, 1)$ such that $\alpha + \beta < 1$, and set 
    \[
    C_n := (\log n)^\alpha, \quad B_n := (\log n)^\beta.
    \]
    Assume that $M(\lambda)$ exists for all $\lambda \in \R$, and that $\max_{i \in [n]}|X_i| \leq B_n$ almost surely. Then for all $\eps > 0$,
    \[
    \sup_{\lambda \in [-C_n, C_n]} |M_n^{(k)}(\lambda) - M^{(k)}(\lambda)|= O_p(n^{-1/2 + \eps}).
    \]
\end{proposition}
\begin{proof}
    First we note that, because $\alpha + \beta < 1$, we have 
    \begin{equation}
        \exp(C_nB_n) = \exp((\log n)^{\alpha + \beta}) = n^{(\log n)^{\alpha + \beta - 1}} = n^{o(1)}.
        \label{eqn:alpha-beta}
    \end{equation}
    For each $n$, let $\Pi_n \subset [-C_n, C_n]$ be a $n^{-1}$-covering set of $[-C_n, C_n]$ of size $N:=C_nn + 1$. Then, for any $\lambda \in [-C_n, C_n]$, there is $\pi_\lambda \in \Pi_n$ such that $|\lambda - \pi_\lambda| < n^{-1}$, and the following holds by the triangle inequality:
    \begin{equation}
        |M_n^{(k)}(\lambda) - M^{(k)}(\lambda)| \leq |M_n^{(k)}(\lambda) - M_n^{(k)}(\pi_\lambda)| + |M_n^{(k)}(\pi_\lambda) - M^{(k)}(\pi_\lambda)| + |M^{(k)}(\pi_\lambda) - M^{(k)}(\lambda)|.\label{eqn:cover-proj-ineq} 
    \end{equation}
    Let us first address the middle term of (\ref{eqn:cover-proj-ineq}). By Proposition~\ref{propn:mgf-concentration-ineq} and then (\ref{eqn:alpha-beta}), 
    \begin{multline*}
\P\left(\max_{\pi \in \Pi_n} |M_n^{(k)}(\pi) - M^{(k)}(\pi)| > n^{-1/2 + \eps}\right) \leq N\P\left(|M_n^{(k)}(\pi) - M^{(k)}(\pi)| > n^{-1/2 + \eps}\right) \\ 
        \leq 2N\left\{\exp\left(\frac{-2n^{1/2}}{4B^2_n}\right) + \exp\left(\frac{-n^{2\eps}}{8B_n^ke^{4C_nB_n + 1}}\right)\right\}\\
        \leq 2\{n(\log n)^\alpha + 1\}\left\{\exp\left(\frac{-2n^{1/2}}{4(\log n)^{2\beta}}\right)+ \exp\left(\frac{-n^{2\eps}}{8(\log n)^{k\beta}n^{o(1)}}\right)\right\} 
    \end{multline*}
    which vanishes in the limit $n \to \infty$. 
   It then follows that
    \begin{equation}
        \max_{\pi \in \Pi_n} |M_n^{(k)}(\pi) - M^{(k)}(\pi)| = O_p(n^{-1/2 + \eps}).
        \label{eqn:emgf-cover-bound}
    \end{equation}
    Next, with probability one, $\max_i |X_i| \leq B_n$. On that event, for all $\lambda \in [-C_n, C_n]$,
    \[
    M_n^{(k+1)}(\lambda) = \frac1n\sum_{i=1}^n(X_i - \bar X)^{k+1}e^{\lambda |X_i - \bar X|} \leq (2B_n)^{k+1}e^{2C_nB_n},
    \]
    which implies by the mean-value theorem that for any $\lambda, \lambda' \in [-C_n, C_n]$, there is $\tilde \lambda$ such that 
    \[
    |M_n^{(k)}(\lambda) - M_n^{(k)}(\lambda')| \leq |M_{n}^{(k+1)}(\tilde \lambda)||\lambda - \lambda'| \leq (2B_n)^{k+1}e^{2C_nB_n}|\lambda - \lambda'|
    \]
    and in particular 
    \begin{equation}
        |M_n^{(k)}(\lambda) - M^{(k)}_n(\pi_\lambda)| \leq (2B_n)^{k+1}e^{2C_nB_n}n^{-1}.\label{eqn:emgf-lipschitz}
    \end{equation}
    Identical reasoning then yields
    \begin{equation}
        |M^{(k)}(\lambda) - M^{(k)}(\pi_\lambda)| \leq (2B_n)^{k+1}e^{C_nB_n}n^{-1}.\label{eqn:mgf-lipschitz}
    \end{equation}
    Now by (\ref{eqn:alpha-beta}), we have 
    \[
    2B_n^{k+1}e^{C_nB_n}n^{-1} = 2(\log n)^{\beta(k+1)} n^{-1+o(1)} = o(n^{-1/2}),
    \]
    so it follows that, with probability one,
    \begin{equation}
        |M_n^{(k)}(\lambda) - M^{(k)}_n(\pi_\lambda)| = o(n^{-1/2}), \quad |M^{(k)}(\lambda) - M^{(k)}(\pi_\lambda)| = o(n^{-1/2}).
        \label{eqn:emgf-final-bound}
    \end{equation}
    We conclude by putting (\ref{eqn:emgf-cover-bound}) and (\ref{eqn:emgf-final-bound}) into (\ref{eqn:cover-proj-ineq}) to get 
    \[
    |M_n(\lambda) - M(\lambda)| \leq o(n^{-1/2}) + O_p(n^{-1/2 + \eps}) + o(n^{-1/2}) = O_p(n^{-1/2 + \eps}).
    \]
\end{proof}
Our next goal is to extend the result of Proposition~\ref{prop:bounded-convergence-lemma} to sub-Gaussian random variables. To do this, we first recall some facts about them.
\begin{proposition}
\label{prop:sub-Gaussian-props}
    Let $X$ be a sub-Gaussian random variable with parameter $\xi_*$. Then: 
    \begin{enumerate}
        \item[(a)] For all $t > 0$, $\P(|X - \E X| > t) \leq \exp\left(-t^2(2\xi_*)^{-2}\right).$
        \item[(b)] For all $k \in\N$, $\E(|X|^{k}) \leq 2\xi_*^k\Gamma(k/2 + 1)$.
        \item[(c)] For all $\lambda \in \R$, $M(\lambda) \leq \exp(2^{-1}\lambda^2\xi_*^2)$.
    \end{enumerate}
\end{proposition}
We shall also need the following generalization of H\"older's inequality: 
\begin{proposition}[\cite{folland1999real}, Exercise 6.31]
\label{prop:generalized-holder}
    Let $1 \leq p, q, r \leq \infty$ satisfy $p^{-1} + q^{-1} + r^{-1} = 1$, and let $(X, \mathcal M, \mu)$ be a measure space. If $f \in L^p(\mu), g \in L^q(\mu), h \in L^r(\mu)$, then $fgh \in L^1(\mu)$. Moreover,
    \[
    \int |fgh| \, d\mu \leq \left(\int |f|^p \, d\mu\right)^{\frac1p} \left(\int |g|^q \, d\mu\right)^{\frac1q} \left(\int |h|^r \, d\mu\right)^{\frac1r}.
    \]
\end{proposition}
The idea is to show that sub-Gaussian random variables can be well-approximated by truncating their range. We may then apply Proposition \ref{prop:bounded-convergence-lemma} to conclude. 
\begin{proof}[Proof of Theorem~\ref{prop:sub-gaussian-convergence-lemma}]
    Fix $\beta \in (\alpha, 1- \alpha)$ and let $B_n := (\log n)^\beta$. Then, set $X_i^{\flat} := X_i\1\{|X_i| \leq B_n\}$. Accordingly, let $\bar X^{\flat} = \frac1n\sum_{i=1}^n X_i^{\flat}$, and 
    \[
    M_n^{\flat}(\lambda) := \frac1n\sum_{i=1}^n e^{\lambda (X^{\flat}_i - \bar X^{\flat})},  \quad M^\flat(\lambda) := \E[e^{\lambda X^{\flat}_1}].
    \]
    Thus, by the triangle inequality, we have 
    \begin{align}
        |M^{(k)}(\lambda) - M_n^{(k)}(\lambda)| \leq &|M^{(k)}(\lambda) - (M^\flat)^{(k)}(\lambda)| \ + \notag\\& |(M^\flat)^{(k)}(\lambda) - (M^\flat_n)^{(k)}(\lambda)| + |(M^\flat_n)^{(k)}(\lambda) - M_n^{(k)}(\lambda)|. 
        \label{eqn:M-flat-ineq}
    \end{align}
    For the first term, we note that if $|X_1| \leq B_n$, then $X^\flat_1 = X_1$. Therefore 
    \begin{equation}
        |M^{(k)}(\lambda) - (M^\flat)^{(k)}(\lambda)| \leq \E\left|(X^\flat_1)^{k}e^{\lambda X^\flat_1} - X_1^{k}e^{\lambda X_1}\right| 
    \leq\E|X_1|^{k}e^{C_nX_1}\1\{|X_1| \geq B_n\}.
        \label{eqn:mgf-cutoff-eq-1}
    \end{equation}
    Then, we use the generalized H\"older inequality (Proposition~\ref{prop:generalized-holder}) with exponents $(3, 3, 3)$, and then apply Proposition~\ref{prop:sub-Gaussian-props} to get 
    \begin{align}
        \E\left|X_1^{k}e^{C_nX_1}\1\{|X_1| \geq B_n\}\right|&\leq\{\E|X_1^{3k}|M(3C_n)\P(X_1 \geq B_n)\}^{1/3} \notag \\
        &\leq \left\{2\Xi^{3k}\Gamma\left(\frac{3k}{2}+1\right)\exp\left(\frac92C_n^2\Xi^2-\frac{B_n^2}{2\Xi^2}\right)\right\}^{1/3}. \label{eqn:mgf-cutoff-eq-2}
    \end{align}
    Now note that since $\alpha < \beta$, 
    \begin{equation}
\exp\left(\frac92C_n^2\Xi^2-\frac{B_n^2}{2\Xi^2}\right) \leq  \exp\left(\frac92\Xi^2(\log n)^{2\alpha}-\frac{(\log n)^{2\beta}}{2\Xi^2}\right)
        \label{eqn:exp-C_n-B_n}
    \end{equation}
    which decays faster than any polynomial; in particular it is $o(n^{-2})$. 
    Thus, putting (\ref{eqn:mgf-cutoff-eq-1}), (\ref{eqn:mgf-cutoff-eq-2}), and (\ref{eqn:exp-C_n-B_n}) together shows that 
    \begin{equation}
        |M^{(k)}(\lambda) - (M^\flat)^{(k)}(\lambda)| = o(n^{-1/2}).
        \label{eqn:mgf-cutoff-eq-3}
    \end{equation}
    The bound for the sample analogue of (\ref{eqn:mgf-cutoff-eq-3}) follows the same strategy, but is slightly more involved. To begin, we note that if, $|X_j| \leq B_n$ for all $j \in [n]$, then $X^\flat_j = X_j$ so  $\bar X^\flat = \bar X$. Therefore 
    \begin{multline}
        \E |(M_n^\flat)^{(k)}(\lambda) - M_n^{(k)}(\lambda)|\leq \frac1n\sum_{i=1}^n \E|(X^\flat_i - \bar X^\flat)^ke^{\lambda (X^\flat_i - \bar X^\flat)} - (X_i - \bar X)^ke^{\lambda (X_i - \bar X)}| \\ = \frac1n\sum_{i=1}^n\E\left|\left\{(X^\flat_i - \bar X^\flat)^ke^{\lambda (X^\flat_i - \bar X^\flat)} - (X_i - \bar X)^ke^{\lambda (X_i - \bar X)}\right\} \1\{\exists j:|X_j| > B_n\}\right|.
        \label{eqn:emgf-cutoff-1}
    \end{multline}
    Then the triangle inequality gives us two terms to bound:
    \begin{multline}
        \frac1n\sum_{i=1}^n\E\left|\left\{(X^\flat_i - \bar X^\flat)^ke^{\lambda (X^\flat_i - \bar X^\flat)} - (X_i - \bar X)^ke^{\lambda (X_i - \bar X)}\right\} \1\{\exists j:|X_j| > B_n\}\right| \\
        \leq \frac1n\sum_{i=1}^n \E\left|(X^\flat_i - \bar X^\flat)^ke^{\lambda (X_i^\flat - \bar X^\flat)}\1\{\exists j: |X_j| > B_n\}\right|\\  + \frac1n\sum_{i=1}^n \E\left|(X_i - \bar X)^ke^{\lambda (X_i - \bar X)}\1\{\exists j: |X_j| > B_n\}\right|
        \label{eqn:emgf-cutoff-2}
    \end{multline}
    The first one can be dealt with by evaluating the expectation, using a union bound, Chernoff's inequality, and then (\ref{eqn:exp-C_n-B_n}):
    \begin{align}
        \frac1n\sum_{i=1}^n \E\left|(X^\flat_i - \bar X^\flat)^ke^{\lambda (X_i^\flat - \bar X^\flat)}\1\{\exists j: |X_j| > B_n\}\right|  
        &\leq \frac1n\sum_{i=1}^n \E\left|(2B_n)^ke^{2C_nB_n}\1\{\exists j: |X_j| > B_n\}\right| \notag \\
        &= \frac1n\sum_{i=1}^n (2B_n)^ke^{2C_nB_n}\P(\exists j: |X_j| > B_n) \notag\\
        &\leq \frac1n\sum_{i=1}^n (2B_n)^ke^{2C_nB_n}n\P(|X_1| > B_n)  \notag \\&\leq(2B_n)^ke^{2C_nB_n}n\exp\left(-\frac{B_n^2}{2\Xi^2}\right) = o(n^{-1/2}).
        \label{eqn:emgf-cutoff-3}
    \end{align}
    Next, we move onto the second term of (\ref{eqn:emgf-cutoff-2}). By the generalized H\"older inequality and then a union bound,
    \begin{align}
        \frac1n\sum_{i=1}^n \ &\E\left|(X_i - \bar X)^ke^{\lambda (X_i - \bar X)}\1\{\exists j: |X_j| > B_n\}\right| \notag \\ &\leq \frac1n\sum_{i=1}^n \{\E(|X_i - \bar X|^{3k})\E[e^{3\lambda(X_1 - \bar X)}]\P(\exists j: |X_j| > B_n)\}^{1/3} \notag\\
        &\leq \frac1n\sum_{i=1}^n \{\E(|X_i - \bar X|^{3k})\E[e^{3\lambda(X_1 - \bar X)}]n\P(|X_1| > B_n)\}^{1/3}
        \label{eqn:emgf-cutoff-4}
    \end{align}
    Now note that $\bar X$ is sub-Gaussian with parameter $n^{-1}\xi_*$, so $X_1 - \bar X$ is sub-Gaussian with parameter $(\xi_*^2 + n^{-1}\xi_*^2)^{1/2} \leq \sqrt 2 \xi_* \leq \sqrt 2 \Xi$. Therefore by Proposition~\ref{prop:sub-Gaussian-props}(c), $\E[e^{3\lambda(X_1 - \bar X)}] \leq \exp(9\lambda^2\Xi^2)$. Then, applying Chernoff's inequality to (\ref{eqn:emgf-cutoff-4}), using the fact that $|\lambda| \leq C_n$, and then finally using (\ref{eqn:exp-C_n-B_n}), we have 
    \begin{multline}
         \frac1n\sum_{i=1}^n \{\E(|X_i - \bar X|^{3k})\E[e^{3\lambda(X_1 - \bar X)}]n\P(|X_1| \geq B_n)\}^{1/3} \\ \leq   \frac1n\sum_{i=1}^n\{\E(|X_i - \bar X|^{3k})\}^{1/3}\exp(3\lambda^2\Xi^2)\exp\left(-\frac{B_n^2}{6\Xi^2}\right) \\ 
 \leq \frac1n\sum_{i=1}^n 2^{k/2}\Xi^k\{\Gamma(3k/2+1)\}^{1/3}\exp\left(3C_n^2\Xi^2 - \frac{B_n^2}{6\Xi^2}\right) 
 = o(n^{-1/2}).
 \label{eqn:emgf-cutoff-5}
    \end{multline}
    Combining (\ref{eqn:emgf-cutoff-1}), (\ref{eqn:emgf-cutoff-2}), (\ref{eqn:emgf-cutoff-3}), (\ref{eqn:emgf-cutoff-4}), and (\ref{eqn:emgf-cutoff-5}) shows that $\E|(M_n^\flat)^{(k)}(\lambda) - M_n^{(k)}(\lambda)| = o(n^{-1/2})$. Then Markov's inequality yields 
    \begin{equation}
        \P(|(M_n^\flat)^{(k)}(\lambda) - M_n^{(k)}(\lambda)| > n^{-1/2})  \leq \frac{2|(M_n^\flat)^{(k)}(\lambda) - M_n^{(k)}(\lambda)|}{n^{-1/2}} = o(1).
        \label{eqn:emgf-cutoff-6}
    \end{equation}
    The parts are in place to conclude. Since the $X^\flat_i$'s are bounded by $(\log n)^\beta$, they satisfy the hypotheses of Proposition~\ref{prop:bounded-convergence-lemma}. Therefore we may apply it to the middle term of (\ref{eqn:M-flat-ineq}). As we have shown before, the first and third terms can be handled by (\ref{eqn:mgf-cutoff-eq-3}) and (\ref{eqn:emgf-cutoff-6}), respectively. Thus, 
    \[
    |M_n^{(k)}(\lambda) - M^{(k)}(\lambda)| \leq O_p(n^{-1/2})+O_p(n^{-1/2 + \eps}) + o(n^{-1/2})=O_p(n^{-1/2 + \eps}),    \]
    which completes the proof.
\end{proof}
\subsubsection{Consequences of Theorem~\ref{prop:sub-gaussian-convergence-lemma}}
We now address the corollaries of Theorem~\ref{prop:sub-gaussian-convergence-lemma}.
\begin{proof}[Proof of Corollary~\ref{cor:cgf-unif-conv}]
    Let $\eps > 0$. According to Jensen's inequality, $M(\lambda) = \E[e^{\lambda X}] \geq e^{\lambda \E X} = 1$, and $M_n(\lambda) = \frac1n \sum_{i=1}^n e^{\lambda (X_i - \bar X)} \geq e^{\lambda \frac1n\sum_{i=1}^n (X_i - \bar X)} = 1$. Then, by the Cauchy-Schwarz inequality and then the sub-Gaussianity of $X$,
    \begin{align}
        M^{(k)}(\lambda) &= \E[X^ke^{\lambda X}] \leq \{\E X^{2k} M(2\lambda)\}^{1/2} \notag\\ &\leq \{\E X^{2k} \exp(2\lambda^2\Xi^2)\}^{1/2} \leq  \{\E X^{2k} \exp(2C_n^2\Xi^2)\}^{1/2}\notag \\
        &= \{\E X^{2k} \exp(2(\log n)^{2\alpha}\Xi^2)\}^{1/2} = o(n^{\eps/2})
        \label{eqn:M-k-ineq}
    \end{align}
    where the last step uses $2\alpha < 1$ so that $(\log n)^{2\alpha - 1} \to 0$. Now by a direct computation we have 
    \[
    \psi^{(k)}(\lambda) = \frac{M^{(k)}}{M} - \sum_{j = 1}^{k-1} {k - 1 \choose j- 1} \psi^{(j)}(\lambda)\frac{M^{(k - j)}(\lambda)}{M(\lambda)},
    \]
    and similarly
    \[
    \psi^{(k)}_n(\lambda) = \frac{M^{(k)}_n(\lambda)}{M_n(\lambda)} - \sum_{j = 1}^{k-1} {k - 1 \choose j- 1} \psi^{(j)}_n(\lambda)\frac{M^{(k - j)}_n(\lambda)}{M_n(\lambda)}.
    \]
    Also, let us observe that, because $\eps > 0$ is arbitrary in Theorem~\ref{prop:sub-gaussian-convergence-lemma}, it holds that 
    \begin{equation}
        \sup_{|\lambda| \leq C_n} |M_n^{(k)}(\lambda) - M^{(k)}(\lambda)| = O_p(n^{-1/2 + \eps/2}).
        \label{eqn:sub-gaussian-convergence-app}
    \end{equation}
    The remainder of the proof goes by induction on $k$. \\
    \textit{Base case}: By the mean-value theorem applied to $x \mapsto \log x$, 
    \[
    |\psi_n(\lambda) - \psi(\lambda)| \leq \tau^{-1}|M_n(\lambda) - M(\lambda)|
    \]
    where $\tau \in [M_n(\lambda) \wedge M(\lambda), M_n(\lambda) \vee M(\lambda)]$. But $M_n(\lambda) \wedge M(\lambda) \geq 1$, so $\tau \geq 1$ and actually \[
    |\psi_n(\lambda) - \psi(\lambda)| \leq |M_n(\lambda) - M(\lambda)|.
    \]
    \textit{Inductive step:} Suppose the claim holds for any $j < k$. Then we have 
    \begin{multline*}
        |\psi^{(k)}(\lambda) - \psi^{(k)}_n(\lambda)| \\ \leq \left|\frac{M^{(k)}_n(\lambda)}{M_n(\lambda)} - \frac{M^{(k)}(\lambda)}{M(\lambda)}\right| + \sum_{j = 1}^{k-1} {k - 1 \choose j- 1} \left|\psi^{(j)}_n(\lambda)\frac{M^{(k - j)}_n(\lambda)}{M_n(\lambda)} - \psi^{(j)}(\lambda)\frac{M^{(k - j)}(\lambda)}{M(\lambda)}\right|.
    \end{multline*}
    For the rest term, since $M(\lambda) \geq 1$ and $M_n(\lambda) \geq 1$, we have  
    \begin{align}
\left|\frac{M^{(k)}_n(\lambda)}{M_n(\lambda)} - \frac{M^{(k)}(\lambda)}{M(\lambda)}\right| &\leq \left|\frac{M_n^{(k)}(\lambda) - M^{(k)}(\lambda)}{M_n(\lambda)}\right| + \left|\frac{M^{(k)}(\lambda)\{M_n(\lambda) - M(\lambda)\}}{M(\lambda)M_n(\lambda)}\right| \notag \\
&\leq \left|M_n^{(k)}(\lambda) - M^{(k)}(\lambda)\right| + \left|M^{(k)}(\lambda)\{M_n(\lambda) - M(\lambda)\}\right|.
\label{eqn:cor1-mgf-triangle}
    \end{align}
    The first term is $O_p(n^{-1/2 + \eps})$ by (\ref{eqn:sub-gaussian-convergence-app}), and the second is $O_p(n^{-1/2 + \eps})$ by (\ref{eqn:M-k-ineq}) and (\ref{eqn:sub-gaussian-convergence-app}) again. Meanwhile, for each term in the summation, we have 
\begin{multline*}
    \left|\psi^{(j)}_n(\lambda)\frac{M^{(k - j)}_n(\lambda)}{M_n(\lambda)} - \psi^{(j)}(\lambda)\frac{M^{(k - j)}(\lambda)}{M(\lambda)}\right| \\ \leq \left|\psi_n^{(j)} - \psi^{(j)}\right|\left|\frac{M^{(k - j)}_n(\lambda)}{M_n(\lambda)}\right| + |\psi^{(j)}|\left|\frac{M^{(k - j)}_n(\lambda)}{M_n(\lambda)} - \frac{M^{(k - j)}(\lambda)}{M(\lambda)}\right|.
\end{multline*}
The first term is $O_p(n^{-1/2 + \eps})$ by the inductive hypothesis and (\ref{eqn:M-k-ineq}). The second follows by similar reasoning to (\ref{eqn:cor1-mgf-triangle}). 
\end{proof}
\begin{proof}[Proof of Proposition~\ref{prop:L_n-cgf-bound}]
    
We claim first that the following formula holds for all $\lambda \in \R$:
    \begin{equation}
        L(\lambda) = 2\int_0^1 (1 - x)\psi''(\lambda x) \, dx, \quad L_n(\lambda) = 2\int_0^1 (1 - x)\psi''_n(\lambda x) \, dx.  \label{eqn:L-integral-rep}
    \end{equation}
    To see this, note that $\psi(0) = \psi'(0) = \E X= 0$ so by Taylor's theorem and then a change of variable,
    \[
    \psi(\lambda) = \int_0^\lambda (\lambda - x)\psi''(x) \, dx = \lambda^2\int_0^1 (1 - x) \psi''(\lambda x) \, dx 
    \]
    so the claim follows by recalling $L(\lambda) = 2\lambda^{-2}\psi(\lambda)$. Meanwhile, since $\psi''(0) = \sigma^2$, we have 
    \[
    L(0) = 2\int_0^1 (1 - x) \psi''(0) \, dx.
    \]
    Identical reasoning holds for $L_n$, noting that $\psi'_n(0) = 0$ too. Now since $\psi^{(k)}$ is twice continuously differentiable, we may differentiate under the integral sign to obtain
    \[
    L^{(k)}(\lambda) = 2\int_0^1 (1 - x)x^k \psi^{(k+2)}(\lambda x) \, dx , \quad L_n^{(k)}(\lambda) = 2 \int_0^1(1 - x)x^k \psi^{(k+2)}_n(\lambda x)\, dx.
    \]
    Then it follows that, for any $\lambda \in [-C_n, C_n]$, 
    \begin{align*}
         |L_n^{(k)}(\lambda) - L^{(k)}(\lambda)|&\leq \int_0^1 (1 - x)x^k|\psi^{(k+2)}(\lambda x) - \psi_n^{(k+2)}(\lambda x)| \, dx \\ &\leq 2\sup_{t \in [-C_n, C_n]} |\psi^{(k+2)}(t) - \psi^{(k+2)}_n(t)| \int_0^1 (1 - x)x^k \, dx \\
        &\leq \sup_{t \in [-C_n, C_n]} |\psi^{(k+2)}(t) - \psi^{(k+2)}_n(t)|
    \end{align*}
    where the last step follows because $\int_0^1 (1-x)x^k \, dx = (k+1)^{-1}(k+2)^{-1} \leq 1/2$.
\end{proof}
\begin{proof}[Proof of Theorem~\ref{prop:second-consistency}]
Define $\tilde \xi_n^2 := \sup_{|\lambda|\leq C_n} L(\lambda)$. We then have 
\begin{align*}
    |\hat \xi^2_n - \xi^2_*| &\leq |\hat \xi^2_n - \tilde \xi^2_n| + \xi_*^2 - \tilde \xi_n^2 \\
    &\leq |\hat \xi^2_n - \tilde \xi^2_n| + \sup_{\lambda \in \R} L(\lambda) - \sup_{\lambda \in [-C_n, C_n]} L(\lambda) \\
    &= |\hat \xi^2_n - \tilde \xi^2_n| + \left\{ \sup_{|\lambda| > C_n} L(\lambda) - \sup_{|\lambda| \leq C_n} L(\lambda)\right\} \vee 0 \\
    &= |\hat \xi^2_n - \tilde \xi_n^2| + \delta(C_n) \vee 0.
\end{align*}
    Since $L$ and $L_n$ are continuous, there is $\tilde \lambda_n \in [-C_n, C_n]$ such that $\tilde \xi_n = L(\tilde \lambda_n)$, and $\lambda^*_n \in [-C_n, C_n]$ such that $\hat \xi_n^2 = L_n(\lambda^*_n)$. Then 
    \[
    \hat \xi_n^2 - \tilde \xi_n^2 \leq L_n(\lambda^*_n) - L(\lambda^*_n) \leq \sup_{\lambda \in [-C_n, C_n]} |L_n(\lambda) - L(\lambda)|
    \]
    and 
    \[
    \hat \xi_n^2 - \tilde \xi_n^2 \geq L_n(\tilde \lambda_n) - L(\tilde \lambda_n) \geq -\sup_{\lambda \in [-C_n, C_n]}|L_n(\lambda) - L(\lambda)| 
    \]
    so that $|\hat \xi_n^2 - \tilde \xi^2_n| \leq \sup_{|\lambda| \leq C_n} |L_n(\lambda) - L(\lambda)|$. Therefore 
    \begin{equation}
        |\hat \xi^2_n - \xi_*^2| \leq \sup_{|\lambda| \leq C_n} |L_n(\lambda) - L(\lambda)| + \delta(C_n) \vee 0.
        \label{eqn:L-delta-bound}
    \end{equation}
    (a) By Proposition~\ref{prop:L_n-cgf-bound}(b), we have $\sup_{|\lambda| \leq C_n} |L_n(\lambda) - L(\lambda)| = o_p(1)$ for any $\eps > 0$. Since $C_n \to \infty$, $\delta(C_n) \to 0$. \\
    (b) Whenever $n$ is large enough that $C_n > C_0$, we have $\delta(C_n) \leq 0$, so the claim follows from (\ref{eqn:L-delta-bound}). \\
    (c) Make $r > 0$ so small that $\sup_{|\lambda| > C_0} L(\lambda) < \sup_{|\lambda| \leq C_0} L(\lambda) - r$. It follows that $\xi_*^2 = L(\lambda^*)$ for some $\lambda^* \in [-C_0, C_0]$. Then, since $L_n(\lambda^*) \xrightarrow{\mathrm{a.s.}} L(\lambda^*)$, if $n$ is large enough, we have on one hand
    \[
    \sup_{|\lambda| \leq C_0} L_n(\lambda) \geq L_n(\lambda^*) \geq L(\lambda^*) - \frac r2
    \]
    and on the other
    \[
    \sup_{C_0 < |\lambda| \leq C_n} L_n(\lambda) \leq \sup_{|\lambda| \leq C_n} |L_n(\lambda) - L(\lambda)| + \sup_{C_0 < |\lambda|} L(\lambda) \leq o_p(1) + \xi_*^2 - r  \leq \xi_*^2 - \frac{r}{2}
    \]
    so that
    \[
    \sup_{|\lambda| \leq C} L_n(\lambda) > \sup_{C < |\lambda| \leq C_n} L_n(\lambda).
    \]
    It follows that none of $(C, C_n]$ can be a maximizer. Thus with high probability, the maximizers of $L_n$ contained in $[-C_0, C_0]$, so it actually holds that 
    \[
    |\hat \xi^2_n - \tilde \xi^2_n| \leq \sup_{|\lambda| \leq C_0} |L_n(\lambda) - L(\lambda)|.
    \]    Proposition~\ref{prop:L_n-cgf-bound}(a) then applies to yield $\sqrt n(\hat \xi_n^2 - \tilde \xi^2_n) = O_p(1)$. Since the hypothesis of (c) clearly implies (b), it holds again that $\delta(C_n) \vee 0 = 0$ for large enough $n$, and the claim follows from (\ref{eqn:L-delta-bound}).
\end{proof}
%{\color{red} Some further remarks about Theorem~\ref{prop:second-consistency} can be made. Because the bound in Theorem~\ref{prop:second-consistency} is valid for finitely many samples, it can also be understood as a result that gives three regimes that $\hat \xi_n^2$ must obey, regardless of the distribution. Let $C_n \to \infty$ be any increasing sequence. 
%\begin{enumerate}
%    \item[(a)] (\textbf{Slow Regime}) For any $C_n$ such that $\delta(C_n) > 0$, we have 
%    \[|\hat \xi_n^2 - \xi_*^2| \leq \sup_{|\lambda| \leq C_n} |L_n(\lambda) - L(\lambda)| + \delta_P(C_n).\]
%    When $C_n = (\log n)^\alpha$ for $\alpha \in (0, 1/2)$, the first term is $O_p(n^{-1/2 + \eps})$ for any $\eps > 0$. Meanwhile, if $\delta_P(C) \lesssim C^{-1}$ (true for bounded random variables), we have $\delta_P(C_n) = O((\log n)^{-\alpha})$ and therefore the $\delta_P(C_n)$ term contributes the largest amount of estimation error. 
 %   \item[(b)] (\textbf{Intermediate Regime}) For any $C_n$ such that $\delta(C_n) \leq 0$, we have 
%    \[|\hat \xi_n^2 - \xi_*^2| \leq \sup_{|\lambda| \leq C_n} |L_n(\lambda) - L(\lambda)|,\]
  %  and thus the regime is strictly faster. If $C_n = (\log n)^\alpha$, then the rate is $O_p(n^{-1/2 + \eps})$ for all $\eps > 0$. 
  %  \item[(c)] (\textbf{Parametric Regime}) Let $n_0$ be the first integer for which $\delta(C_{n_0}) < 0$. Then 
%    \[|\hat \xi_n^2 - \xi_*^2| \leq \sup_{|\lambda| \leq C_{n_0}} |L_n(\lambda) - L(\lambda)|,\]
%    which gives the exact rate of $O_p(n^{-1/2})$. 
%\end{enumerate}
%This formulation shows that taking $\alpha \to 1/2$ gives a faster entry into the intermediate and parametric regimes. Future work on data-driven choices of $C_n$ could then give faster access to the regimes in (b) and (c).

%Also note that for some certain choices $P$ (example: $\frac12 N(0, 1) + \frac12 N(0, 2)$), we have $\delta_P(C_n) > 0$ for all $n$, and thus we remain in the slow regime for all $n$. }
\subsubsection{Proof of Proposition~\ref{prop:ptwise-asymptotic-normality}}
In this section, we show the asymptotic normality of $\hat \xi_n^2$. We shall first need a lemma: 
\begin{proposition}
\label{prop:maximizer-root-n}
    Suppose that $\limsup_{|\lambda| \to \infty} L(\lambda) < \xi_*^2$ and that $\inf \{L''(\lambda^*): \lambda^* \in \Lambda^*\} < 0$. Define $\pi: \R \to \Lambda^*$ by the rule 
    \[
    \pi(\lambda) := \argmin_{\lambda^* \in \Lambda^*} |\lambda - \lambda^*|.
    \]
    Then, for any sequence of maximizers $\lambda^*_n \in \Lambda^*_n$, we have $\lambda^*_n - \pi(\lambda^*_n) = O_p(n^{-1/2})$. 
\end{proposition}
\begin{proof}
    Fix $\eps > 0$. Since $\limsup_{|\lambda| \to \infty} L(\lambda) < \xi_*^2$, we may choose $C, \delta_0 > 0$ such that 
    \[
    \sup_{|\lambda| > C} L(\lambda) < \xi_*^2 - 2\delta_0.
    \]
    Then $\Lambda^* \subset [-C, C]$ is bounded and hence compact, so that  $\pi$ makes sense. Furthermore, 
    the set \[
    A_\eps := [-C, C] - \bigcup_{\lambda^* \in \Lambda^*} \{\lambda: |\lambda-\lambda^*| \leq \eps\}
    \]
    is compact and contains no maximizer of $L$, so $L$ attains some maximum $\xi_A^2 < \xi_*^2$. Put $\delta := (\xi_*^2 - \xi_A^2)/2$, and define $E_n := \{|\lambda^*_n - \lambda^*| > \eps\}$. On $E_n$, if $\lambda^*_n \in [-C, C]$ then $\lambda^*_n \in A_\eps$ so $L(\lambda^*_n) \leq \xi_A^2$. If $\lambda^*_n \not\in [-C, C]$ then $L(\lambda^*_n) \leq \sup_{|\lambda| > C} L(\lambda) \leq \xi_A^2$. Thus we have $L(\lambda^*_n) \leq \xi_A^2$ whenever $E_n$ holds. It then follows that 
    \[
    \hat \xi_n^2 = L_n(\lambda^*_n) \leq L(\lambda^*) + \sup_{|\lambda| \leq C_n} |L_n(\lambda) - L(\lambda)|
    \]
    and 
    \[
    \hat \xi_n^2 = L_n(\lambda^*_n) \geq L_n(\lambda^*) \geq L(\lambda^*) - \sup_{|\lambda| \leq C_n} |L_n(\lambda) - L(\lambda)| 
    \]
    so that 
    \[
    \xi_*^2 - \xi_A^2 \leq 2\sup_{|\lambda| \leq C_n} |L_n(\lambda) - L(\lambda)|.
    \]
We thus have $E_n \subset \{\sup_{|\lambda| \leq C_n} |L_n(\lambda) - L(\lambda)| \geq \delta\}$. The probability of the second event vanishes according to Proposition~\ref{prop:L_n-cgf-bound}(b), so we conclude $\lambda^*_n - \pi(\lambda^*_n) = o_p(1)$.

    We now deduce the rate of convergence.
    By the mean-value theorem, there is $\tilde \lambda_n$ between $\lambda^*_n$ and $\pi(\lambda^*_n)$ such that 
    \[
    0 = L_n'(\lambda^*_n) = L'_n(\pi(\lambda^*_n)) + (\lambda^*_n - \pi(\lambda^*_n))L''_n(\tilde \lambda_n),
    \]
    and by the triangle inequality, we have 
    \[
    |L''_n(\tilde \lambda_n) - L''(\pi(\lambda^*_n))| \leq |L''_n(\tilde \lambda_n) - L''(\tilde \lambda_n)| + |L''(\tilde \lambda_n) - L''(\pi(\lambda^*_n))|. 
    \]
    The first term is $o_p(1)$ by Proposition~\ref{prop:L_n-cgf-bound}(a). By the boundedness of $\Lambda^*$ and the fact that $\tilde \lambda_n - \pi(\lambda^*_n) = o_p(1)$, we may restrict $L''$ to a compact interval containing all of the $\tilde \lambda_n$'s and $\Lambda^*$. On this interval, $L''$ must be uniformly continuous, so it follows that the second term is $o_p(1)$ as well. Thus we have $L''_n(\tilde \lambda_n) - L''(\pi(\lambda^*_n)) = o_p(1)$. Since $\inf \{L''(\lambda^*): \lambda^* \in \Lambda^*\} < 0$,  provided $n$ be large enough, $L''_n(\tilde \lambda^*_n)$ will be bounded away from zero. Therefore, we have 
    \[
    |\pi(\lambda^*_n)-\lambda^*_n|  = \left|\frac{L'_n(\pi(\lambda^*_n))}{L''_n(\tilde \lambda_n)}\right| = \left|\frac{L'_n(\pi(\lambda^*_n)) - L'(\pi(\lambda^*_n))}{L''_n(\tilde \lambda_n)}\right|  \leq \sup_{\lambda^* \in \Lambda^*} \frac{|L'_n(\lambda^*) - L'(\lambda^*)|}{|L''_n(\tilde \lambda_n)|}.
    \]
    Since $\Lambda^*$ is bounded, Proposition~\ref{prop:L_n-cgf-bound}(a) says that the convergence of $L'_n$ is uniform and $O_p(n^{-1/2})$. This proves the claim.  
\end{proof}
We are now ready to prove a central limit theorem for $\hat \xi_n^2$. 
\begin{proof}[Proof of Proposition~\ref{prop:ptwise-asymptotic-normality}]
    Let $\lambda^*$ be the unique maximizer of $L$. Since there is $C_0$ such that $\delta(C_0) < 0$, we have that
    \[
    \sup_{|\lambda|  \geq C_0 \vee \lambda^*} L(\lambda) < \sup_{|\lambda| \leq C_0 \vee \lambda^*} L(\lambda) = \xi_*^2
    \]
    and therefore the hypotheses of Proposition~\ref{prop:maximizer-root-n} hold. 
     Then, by Taylor's theorem there is $\tilde \lambda_n$ between $\lambda_n^*$ and $\lambda^*$ such that 
    \begin{align*}
         L_n(\lambda^*_n) - L_n(\lambda^*) &= (\lambda^*_n - \lambda^*) L'_n(\lambda^*) + \frac{(\lambda_n^* - \lambda^*)^2}{2}L''_n(\tilde \lambda_n) \\
         &= O_p(n^{-1/2})O_p(n^{-1/2}) + O_p(n^{-1})O_p(1) = O_p(n^{-1})
    \end{align*}
    by Proposition~\ref{prop:L_n-cgf-bound}(a). 
    Therefore, it holds that 
    \begin{align*}
        \sqrt n(\hat \xi_n^2 - \xi_*^2) &= \sqrt n(L_n(\lambda^*) - L(\lambda^*))  + \sqrt n(L_n(\lambda^*_n) - L_n(\lambda^*)) \\ &=  \sqrt n(L_n(\lambda^*) - L(\lambda^*)) + O_p(n^{-1/2}) \\
        &= \sqrt n(L_n(\lambda^*) - L(\lambda^*)) + o_p(1).
    \end{align*}
    Now let $\lambda$ be fixed. By the central limit theorem, we have 
    \[
    \sqrt n\begin{pmatrix}
        \tilde M_n(\lambda) - M(\lambda) \\ \lambda\bar X
    \end{pmatrix} \xrightarrow{d} N\left(0, \Sigma\right)
    \]
    where 
    \[
    \Sigma = \begin{pmatrix}
        M(2\lambda) - M(\lambda)^2 & \lambda M'(\lambda) \\ \lambda M'(\lambda) & \lambda^2\sigma^2
    \end{pmatrix}.
    \]
    Now since 
    \[
    L_n(\lambda) = \frac{2}{\lambda^2}\left(\log \tilde M_n(\lambda) - \lambda\bar X\right)
    \]
    we have by the Delta method that 
    \[
    \sqrt n (L_n(\lambda) - L(\lambda)) \xrightarrow{d} N(0, V(\lambda))
    \]
    where 
    \[
    V(\lambda) = \frac{4}{\lambda^4}\begin{pmatrix}
        1/M(\lambda) \\ -1
    \end{pmatrix}^\top \Sigma \begin{pmatrix}
        1/M(\lambda) \\ 1
    \end{pmatrix} = \frac{4}{\lambda^4} \left(\frac{M(2\lambda)}{M(\lambda)^2} - 2\lambda \psi'(\lambda)+  \lambda^2\sigma^2 - 1\right).
    \]
\end{proof}
\subsection{Proofs for Section 3}
In order to simplify the proof of Proposition~\ref{prop:minimax-rates}, we first handle Proposition~\ref{prop:unif-supp-delta-inclusion}.
\begin{proof}[Proof of Proposition~\ref{prop:unif-supp-delta-inclusion}]
    Let $P$ be supported on $[-1, 1]$. Then, for any $\lambda \in \mathbb{R}$, 
    \[
    L(\lambda; P) \leq \frac{2}{\lambda^2} \log e^{ |\lambda|} \leq \frac{2}{|\lambda|}.
    \]

    Combining this with the fact that $L(0 ; P) = \sigma^2(P) \geq 0$, we have that, for any $C > 0$,
    \begin{align*}
    \delta_P(C) = \sup_{|\lambda| \geq C} L(\lambda; P) - \sup_{| \lambda | \leq C} L(\lambda; P) 
    \leq \sup_{|\lambda| \geq C} L(\lambda ; P) \leq \sup_{|\lambda| \geq C} \frac{2}{|\lambda|} \leq \frac{2}{C}.
    \end{align*}
    
   The first claim of the Proposition thus follows. 
   
   
   Now suppose $\sigma^2(P) \geq \frac32\delta_0$. Let $C_1 := 8 \delta_0^{-1}$ and observe that
   \[
   \sup_{|\lambda| \geq C_1} L(\lambda; P) \leq \sup_{|\lambda| \geq C_1} \frac{4}{|\lambda|}  = \frac{4}{C_1} \leq \frac{\delta_0}{2}.
   \]

    Therefore, we have 
    \begin{align*}
    \delta_P(C_1) &= \sup_{|\lambda| \geq C_1} L(\lambda ; P) - \sup_{|\lambda| \leq C_1} L(\lambda ; P) \leq \frac{\delta_0}{2} - L(0 ; P) \\
    &= \frac{\delta_0}{2} - \sigma^2(P) \leq - \delta_0,
    \end{align*}
as desired. The Proposition thus follows. 
\end{proof}

Now we address the minimax lower bounds. In preparation for this, we shall need to grasp the sub-Gaussian parameter of symmetric three-point distributions. In some cases, the sub-Gaussian parameter can be computed explicitly: 
\begin{proposition}[\cite{atouani2025optimalsubgaussianvarianceproxy}, Theorem 3.1(a)]
\label{prop:3-point-sub-gauss}
    For $a, p > 0$, let $P := (1 - 2p)\delta_0 + p\delta_{-a} + p\delta_a$. If $p \in [1/6, 1/2)$, then $\xi_*^2(P) = \sigma^2(P) = 2pa^2$.
\end{proposition}
While \cite{atouani2025optimalsubgaussianvarianceproxy} gives exact expressions for the sub-Gaussian parameter of three-point distributions when $p \leq 1/6$, they are still rather difficult to work with. Instead, it will be easier to settle for an estimate of $\xi_*^2(P)$ via Proposition~\ref{prop:three-point-computation}. To simplify the proof, we first prove that the sub-Gaussian parameter scales like variance: 
\begin{proposition}
\label{prop:subgauss-scaling}
    Let $X \sim P$ be sub-Gaussian. If $aX \sim P_a$, then $\xi_*^2(P) = a^2\xi_*^2(P_a)$.
\end{proposition}
\begin{proof}
    That $\xi_*^2(P) \leq a^2\xi_*^2(P_a)$ is clear from the definition. For the lower bound, let $(\lambda_n)$ be a sequence such that  $L(\lambda_n; P) \to \sup_\lambda L(\lambda; P)$. It then follows that $L(\lambda_n/a; P_a) = a^2L(\lambda_n; P)$, so $ L(\lambda_n/a; P_a) \to a^2\xi_*^2(P)$, and thus $\xi_*^2(P_a) \geq a^2 \xi_*^2(P)$.
\end{proof}
\begin{proposition}
    
\label{prop:three-point-computation}
    For a real sequence $(a_n)$, let $P := (1 - \frac1n)\delta_0 + \frac{1}{2n}\delta_{a_n} + \frac{1}{2n} \delta_{-a_n}$. The following hold:
    \begin{enumerate}
        \item[(a)] $\frac{a_n^2}{2\log 2n} \leq \xi_*^2(P) \leq \frac{a_n^2}{\log n}$. 
        %\item[(b)] Let $\lambda^*$ be a $[0, \infty)$-maximizer of $L(\lambda; P)$; we have $\lambda^* \leq 2a_n^{-1}\log 2n$. 
        \item[(b)] Put $a_n := (\log n)^\gamma$ for $\gamma \in [0, 1/2]$. Then $\delta_P(C) \leq  r_\gamma(C)$ for all $C > 0$.
    \end{enumerate}
\end{proposition}
\begin{proof}
    (a) Because we can rescale by $a_n$ and then apply Proposition~\ref{prop:subgauss-scaling}, it suffices to prove that if  $P = (1 - \frac1n)\delta_0 + \frac{1}{2n}\delta_{1} + \frac{1}{2n} \delta_{-1}$ then 
    \[
     \frac{1}{2 \log 2n}\leq \xi_*^2(P) \leq \frac{1}{\log n}.
    \]
    First we address the lower bound.  We have \begin{align*}
        M(2 \log 2n; P) &= \left(1 - \frac1n\right) + \frac{1}{2n}e^{2 \log 2n} + \frac{1}{2n}e^{-2 \log 2n} \\
        &= \left(1 - \frac1n\right) + \frac1n\left(\frac{(2n)^2 +(2n)^{-2}}{2}\right) \\
        &\geq \left(1 - \frac1n\right)  + 2n \geq 2n,
    \end{align*}
    so it follows that 
    \[
    \xi_*^2(P) \geq L(2 \log 2n; P) \geq \frac{2}{(2 \log 2n)^2}\log 2n \geq \frac{1}{2\log 2n}.
    \]
    For the upper bound, we will use the definition of $\xi_*^2$ and show that 
    \begin{equation}
        M(\lambda; P) \leq \exp\left(\frac{\lambda^2}{2\log n}\right).
    \label{eqn:mgf-upper-bound}
    \end{equation}
    There are two cases to handle. First, if $\lambda \leq \sqrt{2\log n}$, we have that 
    \[
    \log M(\lambda) \leq \frac{1}{n}\left(\frac{e^{\lambda} - e^{-\lambda}}{2} - 1\right) \leq\frac{1}{n}\left(e^{\lambda^2/2} - 1\right)
    \]
    by the inequalities $\log(1+x)\leq x$ and $\frac{e^x - e^{-x}}{2} \leq e^{x^2/2}$. Then since $s \mapsto (e^s - 1)/s$ is increasing and $\lambda^2/2 \leq \log n$ we have 
    \[
\frac{e^{\lambda^2/2}-1}{\lambda^2/2} \leq  \frac{e^{\log n} - 1}{\log n} \leq \frac{n}{\log n}
    \]
    so multiplying by $\lambda^2/2$ yields $e^{\lambda^2/2} - 1 \leq \frac{n\lambda^2}{2\log n}$. Therefore
    \[
    \log M(\lambda) \leq \frac{\lambda^2}{2 \log n},
    \]
    so $\lambda$ satisfies (\ref{eqn:mgf-upper-bound}) in this case. On the other hand, if $\lambda \geq \sqrt{2 \log n}$, then we have 
    \[
    \log M(\lambda) \leq \frac{1}{n}\left(\frac{e^{\lambda} - e^{-\lambda}}{2} - 1\right)  \leq \frac{e^{\lambda} - 1}{n}.
    \]
    By the AM-GM inequality, $\frac12(\log n + \frac{\lambda^2}{\log n})\geq \lambda$, so it follows that $e^{\lambda } \leq \exp(\frac12(\log n + \frac{\lambda^2}{\log n})) = \sqrt n \exp(\frac{\lambda^2}{2\log n})$. Combining this with the above display, it only remains to prove that $\sqrt n\exp(\frac{\lambda^2}{2\log n}) - 1  \leq n\exp(\frac{\lambda^2}{2\log n})$, but this holds trivially because $\frac{\lambda^2}{2\log n} > 1$. 
    \\
    %(b) Write $\lambda_n = 2a_n^{-1}\log 2n$. By a direct computation, we obtain that $L'(\lambda) > 0$ for all $\lambda \in (0, \lambda^*)$ and $L'(\lambda) < 0$ for all $\lambda \in (\lambda^*, \infty)$. It thus suffices to prove that $L'(\lambda_n) < 0$. Indeed, we have 
    %\[
    %L'(\lambda) = \frac{2}{\lambda^3}[\lambda\psi(\lambda) - 2 \psi(\lambda)],
    %\]
    %and by a direct computation, we have 
    %\[\frac{\lambda_nM'(\lambda_n)}{M(\lambda_n)} =  \lambda_na_n\frac{(2n - 1/8n^3)}{2n + 1 - 1/n + 1/8n^3} \leq \lambda_na_n =2\log 2n.
    %\]
    %Furthermore since $M(\lambda_n) \geq 2n$, $2\psi(\lambda_n) \geq 2\log 2n$. Putting these two inequalities together shows that $\lambda_n\psi(\lambda_n) < 2\psi(\lambda_n)$, and therefore $L'(\lambda_n) < 0$.
    %\\
    (b) Let $C > 0$ be arbitrary. If $2(\log n)^{1 - \gamma} \leq C$ then we have 
    \begin{align*}
        \delta_P(C) &\leq \sup_{|\lambda| \geq C} L(\lambda) \leq \sup_{|\lambda| \geq C} \frac{2}{\lambda^2} \log \E[e^{a_n\lambda}] \\ &\leq \sup_{|\lambda| \geq C}  \frac{2a_n}{\lambda} = \frac{2a_n}{C} \leq r_\gamma(C)
    \end{align*}
    where the last inequality follows by direct computation because $\gamma \geq 0$. 
    On the other hand, if $C \leq 2(\log n)^{1 - \gamma}$ then by (a), 
    \[
    \delta_P(C) \leq \xi_*^2(P) \leq (\log n)^{2\gamma - 1} = 2^{\frac{1 - 2\gamma}{1 - \gamma}}2^{\frac{2\gamma - 1}{1 - \gamma}} \frac{(\log n)^{\frac{2\gamma - 1}{1 - \gamma}}}{(\log n)^{\gamma\frac{2\gamma - 1}{1 - \gamma}}}=r_\gamma(2(\log n)^{1 - \gamma}) \leq r_\gamma(C).
    \]
\end{proof}
Now we can prove Proposition~\ref{prop:minimax-rates}.
\begin{proof}[Proof of Proposition~\ref{prop:minimax-rates}]
(a) By Proposition~\ref{prop:subgauss-scaling}, it will suffice to prove the case when $\Xi = 1$. In the general case, we rescale all distributions involved by $\sqrt{\Xi}$ so that their sub-Gaussian parameters are rescaled by $\Xi$. Towards this direction, let $a_n =(\log n)^\gamma$. Let $P_1 = \delta_0$, and for $n \in \N$, let $P_2 = (1 - \frac{1}{n})\delta_0 + \frac{1}{2n}\delta_{ a_n} + \frac{1}{2n}\delta_{-a_n}$. Note then that by Propositions~\ref{prop:subgauss-scaling} and \ref{prop:three-point-computation}(a) we have 
\[
\xi_*^2(P_2) = \xi_*^2\left((1 - \frac{1}{n})\delta_0 + \frac{1}{2n}\delta_{(\log n)^\gamma} + \frac{1}{2n}\delta_{-(\log n)^\gamma}\right) =  (\log n)^{2\gamma - 1}\leq 1,
\]
and thus Proposition~\ref{prop:three-point-computation}(b) shows that $\xi_*^2(P_2) \in \mathcal P_c(1, 0, r_\gamma)$.
Now we have 
\begin{align*}
    H^2(P_1, P_2) &= \frac12 \sum_{x \in \{0, \pm a_n\}} \left(\sqrt{P_1(x)} - \sqrt{P_2(x)}\right)^2 \\ &= 1 - \sum_{x \in \{0, \pm a_n\}} \sqrt{P_1(x)P_2(x)} \\
    &= 1 - \sqrt{1 - n^{-1}}.
\end{align*}
Since $1 - H^2(P_1^{\otimes n}, P_2^{\otimes n}) = (1 - H^2(P_1, P_2))^n$, we have 
\begin{align*}
    H^2(P_1^{\otimes n}, P_2^{\otimes n})  = 1 - \left( 1- n^{-1}\right)^{n/2},
\end{align*}
and then we use $\mathrm{TV}(P^{\otimes n}_1, P^{\otimes n}_2) \leq \sqrt 2 H(P^{\otimes n}_1, P^{\otimes n}_2)$ to get 
\begin{align*}
    \mathrm{TV}(P^{\otimes n}_1, P^{\otimes n}_2) \leq \sqrt{2\left(1 - (1 - n^{-1})^{n/2}\right)}
\end{align*}
so that $\limsup_{n \to \infty} \mathrm{TV}(P^{\otimes n}_1, P^{\otimes n}_2) \leq \sqrt{2(1 - e^{-1/2})}$. By Proposition~\ref{prop:three-point-computation}(a), we have $\xi_*^2(P_2) \geq \frac{(\log n)^{2\gamma}}{2 \log 2n}$, and clearly $\xi_*^2(P_1) = 0$. Thus we may apply Le Cam's two point method to obtain that
    \[
    \inf_{\hat \xi^2} \sup_{P \in \mathcal P(1, 0, r_\gamma)} \E_P|\hat \xi^2 - \xi_*^2| \geq \frac{\left(1 - \sqrt{2(1 - e^{-1/2})}\right)}{8}\frac{(\log n)^{2\gamma}}{\log 2n}.
    \] 
    (b) Without loss of generality, we may assume that $C_0 = 16$ and $\delta_0 = 1/2$ since the class $\mathcal{P}(\Xi, C_0, r_0)$ with $r_0(t) = -\delta_0$ becomes larger by increasing $C_0$ or decreasing $\delta_0$. Moreover by the same rescaling argument as in (a), we may assume $\Xi = 1$.
    
    For $q \in (0, 1)$, define $P_q := (1 - q)\delta_0 + (q/2)\delta_1 + (q/2)\delta_{-1}$. Define $q_0 = \frac{1}{2}$ and $q_n := q_0 + 1/(4\sqrt n)$. 
    
    It is clear that $\xi^2_*(P_{q_0})$ and $\xi^2_*(P_{q_n})$ are both no larger than 1 since the distributions are supported on $[-1, 1]$. Moreover, by Proposition~\ref{prop:unif-supp-delta-inclusion}, $P_{q_0}, P_{q_n} \in \mathcal P(\Xi, C_0, r_-)$. Then, by Proposition~\ref{prop:3-point-sub-gauss} and the fact that $q_0$ and $q_n$ are both greater than $1/3$, we have
    \[
    |\xi_*^2(P_{q_0}) - \xi_*^2(P_{q_n})| = |\sigma^2(P_{q_0}) - \sigma^2(P_{q_n})| =  1/(4\sqrt n).
    \]
    On the other hand, we have by Theorem 5 of \cite{gibbs2002choosingboundingprobabilitymetrics} that 
    \begin{align*}
        \text{KL}(P_{q_0}, P_{q_n}) &= q_0 \log \frac{q_0}{q_n} + (1-q_0)\log \frac{1-q_0}{1-q_n} \\
        &= \frac12 \log \frac{1}{2q_n} + \frac12 \log \frac{1}{2(1 - q_n)} \\ &= \text{KL}(\text{Ber}(q_0), \text{Ber}(q_n)) \\ 
        &\leq \chi^2(\text{Ber}(q_0), \text{Ber}(q_n)) \\
        &= \frac{(q_0 - q_n)^2}{4 q_n} + \frac{(q_0 - q_n)^2}{(1- q_n)} \\
        &\leq \frac{(1/4)^2}{n} \biggl( \frac{1}{4q_n} + \frac{1}{1-q_n}\biggr) = \frac{(1/4)^2}{n}\frac{ (1 + 3q_n)}{4 q_n(1 - q_n)} \leq \frac{1}{3n},
    \end{align*}
    so $\text{KL}(P_{q_0}^{\otimes n}, P_{q_n}^{\otimes n} ) = n\text{KL}(P_{q_0}, P_{q_n}) \leq \frac{1}{3}$. By Pinsker's inequality, we have that $\mathrm{TV}(P_{q_0}^{\otimes n}, P_{q_n}^{\otimes n}) \leq \sqrt{ \frac{1}{2} \mathrm{KL}( P_{q_0}^{\otimes n}, P_{q_n}^{\otimes n} )} \leq \sqrt{1/6}$. 
    \[
    \inf_{\hat \xi^2} \sup_{P \in \mathcal P(\Xi, C_0, r_-)} \mathbb E_P|\hat \xi^2(X_1, \dots, X_n) - \xi_*^2(P)| \geq \frac{1}{8\sqrt n}\biggl(1 - \frac{1}{\sqrt{6}}\biggr). 
    \]
\end{proof}
\begin{proof}[Proof of Proposition~\ref{prop:P(Xi,K)-unif-consistency}]
By the triangle inequality and then Theorem~\ref{prop:second-consistency},  
\[
|\hat \xi^2_n - \xi_*^2(P)| \leq |\tilde \xi^2_n(P) - \xi_*^2(P)| + |\hat \xi^2_n - \tilde \xi^2_n(P)|  = |\tilde \xi^2_n(P) - \xi_*^2(P)| + O_p(n^{-1/2 + \eps})
\]
for any $\eps > 0$, uniformly over $\mathcal P(\Xi)$, and 
\[
\xi_*^2(P) - \tilde \xi^2_n(P) \leq \delta_P(C_n) \vee 0 \leq r(C_n)
\]
so $\hat \xi^2_n - \xi_*^2(P) = r(C_n) + O_p(n^{-1/2 + \eps})$.
\end{proof}
\begin{proof}[Proof of Proposition~\ref{prop:maximizer-uniform-bound}]
By assumption, we have $\delta_P(C_0) < -\delta_0$, so 
\[
\sup_{|\lambda| \geq C_0} L(\lambda) \leq \sup_{|\lambda| \leq C_0} L(\lambda) - \delta_0.
\]
Therefore, none of the maximizers of $L$ can be outside of the interval $[-C_0, C_0]$. 

As for its empirical analogue $L_n$, first let $\lambda^* \in \Lambda^*(P)$. By Proposition~\ref{prop:L_n-cgf-bound}(a), there exists a sequence $\eta_n = o(1)$ dependent only on $\Xi$ and $\delta_0$ and $C_0$ such that, with probability at least $1 - \eta_n$, 
    \begin{equation}
        \sup_{|\lambda| \leq C_0} L_n(\lambda) \geq L_n(\lambda^*) \geq L(\lambda^*) - \sup_{|\lambda| \leq C_0} |L_n(\lambda) - L(\lambda)| > \xi_*^2 - \frac{\delta_0}{2}.
        \label{eqn:L_n-trunc-bound}
    \end{equation}
    By Proposition~\ref{prop:L_n-cgf-bound}(b), there exists another sequence $\eta'_n = o(1)$ dependent only on $\Xi$ and $\delta_0$ such that, with probability at least $1 - \eta'_n$, 
    \begin{align}
    \sup_{|\lambda| \leq C_n} |L_n(\lambda) - L(\lambda)| < \frac{\delta_0}{2}.  \label{eq:Lambda_good_event2}
    \end{align}
    On event~\eqref{eq:Lambda_good_event2}, we have that
    \[
    \sup_{C_0 < |\lambda| \leq C_n} L_n(\lambda) \leq \sup_{|\lambda| \leq C_n} |L_n(\lambda) - L(\lambda)| + \sup_{|\lambda| > C_0} L(\lambda) \leq \frac{\delta_0}{2} + \xi_*^2 - \delta_0 \leq \xi_*^2 - \frac{\delta_0}{2}.
    \]

    On intersection of the event~\eqref{eqn:L_n-trunc-bound} and~\eqref{eq:Lambda_good_event2} therefore, we have that
    \[
    \sup_{|\lambda| \leq C_0} L_n(\lambda) > \sup_{C_0 < |\lambda| \leq C_n} L_n(\lambda),
    \]
    so that $\Lambda^*_n \subset [-C_0, C_0]$. We thus have that
    \[
\sup_{P \in \mathcal{P}(\Xi, C_0, r_0) }\mathbb{P}_P(\Lambda^*_n \not\subset [-C_0, C_0]) \leq \eta_n + \eta'_n,
    \]
    so that the Proposition follows as desired. 
\end{proof}

\subsubsection{Uniform Donsker theorems}
\label{sec:uniform-donsker}
The following definitions are due to \cite{vdv&w2012weak}. Let $\mathbb D$ be a normed vector space, and let $\mathcal B(\mathbb D)$ denote the Borel $\sigma$-algebra on $\mathbb D$ (generated by the open sets under the norm topology of $\mathbb D$). For a probability space $(\Omega, \mathscr F, P)$, we define the \textbf{outer expectation} for any function $X: \Omega \to \mathbb R$ (even non-measurable) as 
\[
\E^*_P X := \inf \left\{\E_PY: \text{$Y$ is $(\mathscr F, \mathcal B(\R))$-measurable and $Y \geq X$}\right\}.
\]
If $(\Omega_n, \mathscr F_n, P_n)$ are a sequence of probability spaces and $X_n : \Omega_n \to \mathbb D$ are arbitrary, then $X_n$ \textbf{converges weakly} to a Borel probability measure $P$ on $\mathbb D$ if as $n \to \infty$, 
\[
\E^*_{P_n} f(X_n) \to \E_P f(X)
\]
for any $f: \mathbb D \to \R$ continuous and bounded. If $\mathbb D$ is separable, then by Theorem 1.12.1 of \cite{vdv&w2012weak}, weak convergence is equivalent to 
\[
\lim_{n \to \infty} \sup_{h \in \mathrm{BL}_1(\mathbb D)} |\E^*_Ph(X_n) - \E_Ph(X)| \to 0,
\]
where 
\[
\mathrm{BL}_1(\mathbb D) = \{h: \mathbb D \to \R: \text{$h$ is 1-Lipschitz and $\|h\|_{\mathbb D} < \infty$}\}.
\]
Let $\mathcal F$ be a family of functions on $\R$. For any function $z: \mathcal F \to \R$, we define 
\[
\|z\|_{\mathcal F} := \sup_{f \in \mathcal F} |z(f)|
\]
and make 
\[
\ell^\infty(\mathcal F) := \left\{z: \mathcal F \to \R : \|z\|_{\mathcal F} <\infty\right\}
\]
into a Banach space by endowing it with the norm $\| \cdot \|_{\mathcal F}$. 

 Let $X_1, \dots, X_n$ be a random sample. The $X_i$'s induce an empirical measure $\mathbb P_n := \frac1n \sum_{i=1}^n \delta_{X_i}$.
 We say that $\mathcal F$ is \textbf{$P$-Donsker} if the empirical process if the empirical process $\mathbb G_{n, P} := \sqrt n(\mathbb P_n - P)$ converges weakly in $\ell^\infty (\mathcal F)$ to a tight Borel measurable element $\mathbb G_P$ of $\ell^\infty(\mathcal F)$. That is, 
\begin{equation}
    \lim_{n \to \infty} \sup_{h \in \text{BL}_1} |\E^*_Ph(\mathbb G_{n, P}) - \E h(\mathbb G_P)| = 0.
    \label{eqn:ptwise-donsker}
\end{equation}
If this convergence is uniform over $\mathcal P$, then $\mathcal F$ is \textbf{Donsker uniformly over $\mathcal P$}:
\[
\lim_{n \to \infty} \sup_{P \in \mathcal P} \sup_{h \in \text{BL}_1} |\E^*_Ph(\mathbb G_{n, P}) - \E h(\mathbb G_P)| = 0.
\]
If the mean-zero Gaussian process $\mathbb G_P$ with covariance function $(f, g) \mapsto \E_{P}f(X)g(X) - \E_Pf(X)\E_Pg(X)$ has a version that is a tight Borel measurable element of $\ell^\infty(\mathcal F)$, then $\mathcal F$ is \textbf{$P$-pre-Gaussian.} If $\mathcal F$ is $P$-pre-Gaussian for each $P \in \mathcal P$, and it holds that
\[
\sup_{P \in \mathcal P} \E_P \|\mathbb G_P\|_{\mathcal F} < \infty \quad \text{and} \quad \lim_{\delta \to 0} \sup_{P \in \mathcal P} \E_P \sup_{\rho_P(f-g) < \delta} |\mathbb G_P(f) - \mathbb G_P(g)| = 0
\]
where $\rho_P(f) = \left(\E_P\{f(X) - \E_Pf(X)\}^2\right)^{1/2}$, then $\mathcal F$ is \textbf{pre-Gaussian uniformly over $\mathcal P$}. 

In order to prove Proposition~\ref{prop:mgf-uniform-donsker}, we shall need the following result: 

\begin{proposition}[\cite{vdv&w2012weak}, Theorem 2.8.4]
\label{prop:uniform-donkser}
    Let $\mathcal P$ be a class of distributions, and $\mathcal F$ be a class of measurable functions with envelope function $F$ that satisfies 
    \[
    \lim_{R \to \infty} \sup_{P \in \mathcal P} \E_P(F^2(X)\1\{F(X) > R\}) = 0.
    \]
    Furthermore, assume that 
    \[
    \int_0^\infty \sup_{P \in \mathcal P} \sqrt{\log N_{[\,]}(\eps \|F\|_{P, 2}, \mathcal F, L_2(P))} \, d\eps < \infty
    \]
    where $N_{[\,]}$ denotes bracketing number. Then $\mathcal F$ is Donsker and pre-Gaussian, both uniformly over $\mathcal P$.
\end{proposition}
\begin{proof}[Proof of Proposition~\ref{prop:mgf-uniform-donsker}]
    An envelope function for $\mathcal F_k$ is $F(x) := |x|^ke^{C|x|}$. Note that, by the Cauchy-Schwarz inequality and then Proposition~\ref{prop:sub-Gaussian-props},
    \[
    \sup_{P \in \mathcal P} \E_PF(X) \leq \sup_{P \in \mathcal P}\{\E_P X^{2k} \E_P e^{2CX}\}^{1/2} \leq 2\Xi^{2k}k! \exp(4C^2\Xi^2) < \infty.
    \]
    Then, by the generalized H\"older inequality, Markov's inequality, and then Proposition~\ref{prop:sub-Gaussian-props}, 
    \begin{align*}
        \E_P(F^2(X) \1\{F(X) > R\}) &= \E_P(X^{2k}e^{CX}\1\{F(X) > R\}) \\
        &\leq \{\E_P(X^{6k})M(3C)\P(F(X) > R)\}^{1/3} \\ 
        &\leq R^{-1/3} \{\E_P(X^{6k})M(3C)\E_P F(X)\}^{1/3} \\
        &\leq R^{-1/3}\{8\Xi^{6k}(3k)!\exp(9 C^2\Xi^2)\} \to 0,
    \end{align*}
    as $R \to \infty$. Therefore the envelope condition holds. Moving onto the bracketing entropy, let $P \in \mathcal P(\Xi)$ and $\eps > 0$. Now we have, for any $\lambda, \lambda' \in [-C, C]$,
\[
|x^ke^{\lambda x} - x^ke^{\lambda'x}| \leq |x^{k+1}e^{\tilde \lambda x}| |\lambda - \lambda'| \leq |x^{k+1}e^{Cx}||\lambda - \lambda'|,
\]
where $\tilde \lambda$ is prescribed by the mean-value theorem. Then, by the Cauchy-Schwarz inequality, we have 
\begin{multline}
    \left\|x^{k+1}e^{Cx}\right\|_{P, 2}\leq \{\E_P X^{2k+2}e^{2CX}\}^{1/2} \leq  \left\{\E_PX^{4k+4} M(4C)\right\}^{1/4} \\ \leq \{2\Xi^{4k+4}\Gamma(2k+3)\exp(8C^2\Xi^2)\}^{1/4} =: D(k, C, \Xi)
\end{multline}
so 
\[
\left\|m_\lambda^{(k)} - m^{(k)}_{\lambda'}\right\|_{P, 2} \leq D(k, C, \Xi)|\lambda - \lambda'|.
\]
Set $\delta := \eps/D(k, C, \Xi)$, and let $\Pi_\delta$ be a $\delta$-net over $[-C, C]$ such that $|\Pi_\delta| \leq \lceil2C\delta^{-1}\rceil = \lceil2CD(k, C, \Xi)\eps^{-1}\rceil$. Now for each $\pi \in \Pi_\delta$, we define 
\[
\ell_\pi (x) := m^{(k)}_\pi(x) - \delta\,|x|^{k+1}e^{C|x|} , \quad u_{\pi}(x) := m^{(k)}_\pi(x) + \delta\,|x|^{k+1}e^{C|x|}.
\]
Then, whenever $\lambda \in [\pi - \delta, \pi + \delta]$, we have $m_\lambda^{(k)} \in [\ell_\pi, u_\pi]$, and furthermore 
\[
\left\|u_\pi - \ell_\pi\right\|_{P, 2} \leq 2\delta D(k, C, \Xi) \leq 2\eps.
\]
We thus have 
\[
\sup_{P \in \mathcal P}N_{[\,]}\!\left(\eps,\, \mathcal F_k,\, L_2(P)\right) 
    \leq
    \left\lceil \frac{4C D(k,C,\Xi)}{\eps} \right\rceil.
\]
Finally, if $\eps > 4CD(k, C, \Xi)$, then $N_{[\,]}\!\left(\eps,\, \mathcal F_k,\, L_2(P)\right)  = 1$, so we have 
\[
\int_0^\infty \sup_{P \in \mathcal P} \sqrt{\log N_{[\,]}\!\left(\eps,\, \mathcal F_k,\, L_2(P)\right)} \leq \int_0^{4CD(k, C, \Xi)} \sqrt{\log \left\lceil \frac{4C D(k,C,\Xi)}{\eps} \right\rceil}\, d\eps < \infty.
\]
We conclude via Proposition~\ref{prop:uniform-donkser}.
\end{proof}


A small technical modification must be made to transfer this property to $M_n(\lambda)$.
\begin{proposition}
\label{prop:centered-mgf-uniform-donsker}
Fix $k \in \N_0$. As $n \to \infty$, \[
\sup_{\lambda \in K} \sqrt n(M_n^{(k)}(\lambda) - M^{(k)}(\lambda; P)) = O_p(1),
\]
uniformly over $P \in \mathcal P(\Xi)$.
\end{proposition}
\begin{proof}
    Recalling that $\tilde M_n(\lambda) := \frac1n\sum_{i=1}^n e^{\lambda X_i}$, we have by Proposition~\ref{prop:mgf-uniform-donsker} that 
    \begin{align*}
        |M_n^{(k)}(\lambda) - M^{(k)}(\lambda)| &\leq |M_n^{(k)}(\lambda) - \tilde M_n^{(k)}(\lambda)| + |\tilde M_n^{(k)}(\lambda) - M^{(k)}(\lambda)| \\ &= |M_n^{(k)}(\lambda) - \tilde M_n^{(k)}(\lambda)| + O_p(n^{-1/2})
    \end{align*}
    uniformly over $P \in \mathcal P(K)$. To address the first term, observe that $M_n(\lambda) = e^{-\lambda \bar X} \tilde M_n(\lambda)$ so by direct differentiation, we have 
    \[
    \sqrt n(M_n^{(k)}(\lambda) - \tilde M_n^{(k)}(\lambda)) = (e^{-\lambda \bar X} - 1)\sqrt n\tilde M_n^{(k)}(\lambda) + e^{-\lambda \bar X} \sum_{j=1}^k {k \choose j}\sqrt n(-\bar X)^j\tilde M_n^{(k - j)}(\lambda).
    \]
    Now note that 
    \[
    \sup_{P \in \mathcal P(\Xi)}\P_P(\sqrt n|\bar X| > t) \leq \sup_{P \in \mathcal P(\Xi)}2\exp\left(-\frac{t^2}{2\xi_*^2}\right) \leq 2\exp\left(-\frac{t^2}{2\Xi^2}\right),\]
    so that $\bar X = O_p(n^{-1/2})$ uniformly over $\mathcal P(\Xi)$. This with Proposition~\ref{prop:mgf-uniform-donsker} shows that all the terms in the summation are $O_p(1)$ uniformly over $\mathcal P(\Xi)$, and therefore we conclude that 
    \[
     \sup_{\lambda \in K}\sqrt n(M_n^{(k)}(\lambda) - \tilde M_n^{(k)}(\lambda)) = O_p(1),
    \]
    uniformly over $\mathcal P(\Xi)$.
\end{proof}
\begin{proof}[Proof of Proposition~\ref{prop:L_n-uniform-rootn}]
    By Proposition~\ref{prop:centered-mgf-uniform-donsker}, we may repeat the argument of Corollary~\ref{cor:cgf-unif-conv} to yield 
    \[
    \sup_{\lambda \in K} \sqrt n|\psi_n^{(k)}(\lambda) - \psi^{(k)}(\lambda)| = O_p(1)
    \]
    uniformly over $\mathcal P(\Xi)$. 
    The claim follows by putting this with Proposition~\ref{prop:L_n-cgf-bound}. 
\end{proof}
\begin{proof}[Proof of Theorem~\ref{prop:hat xi-minimax-optimality}]
    According to Proposition~\ref{prop:maximizer-uniform-bound}, we may make $n$ large enough that the event 
    \[
    \Lambda^*_n \cup \Lambda^*(P) \subset [-C_0, C_0]
    \]
    occurs with high probability, uniformly over $\mathcal P(\Xi, C_0, r_-)$. Moreover, if $C_n$ is made larger than $C$, then we must have $\tilde \xi_n^2(P) = \xi^2_*(P)$ for all $P \in \mathcal P(\Xi, C_0, r_-)$. Therefore we have
    \[
    |\hat \xi_n^2 - \xi_*^2| \leq |\hat \xi^2_n - \tilde \xi^2_n(P)| + |\tilde \xi^2_n(P) - \xi_*^2(P)| \leq |\hat \xi^2_n - \tilde \xi^2_n(P)|.
    \]
    Now repeating the argument in Theorem~\ref{prop:second-consistency}, we have 
    \[
    |\hat \xi^2_n - \tilde \xi^2_n(P)| \leq \sup_{|\lambda| \leq C_0} |L_n(\lambda) - L(\lambda)| = O_p(n^{-1/2}),
    \]
    and this convergence is uniform over $\mathcal P(\Xi, C_0, r_-) \subset \mathcal P(\Xi)$. 
\end{proof}
\subsection{Proofs for Section 4}
Before we present the proof of Proposition~\ref{prop:divergence}, we shall need a brief lemma. 
\begin{proposition}
\label{prop:std-divergence}
    Let $X_1, \dots, X_n \sim P$ and write $\hat \sigma_n^2 = \frac1n \sum_{i=1}^n (X_i - \bar X)^2$. If $\E_PX = \infty$, then $\hat \sigma_n^2 \xrightarrow{\mathrm{a.s.}} \infty$.
\end{proposition}
\begin{proof}
    Observe that 
    \[
    \hat \sigma_n^2 = \frac1n\sum_{i=1}^n (X_i - \bar X)^2 = \frac{1}{2n^2} \sum_{i,j=1}^n (X_i - X_j)^2.
    \]
    Fix $R > 0$, and let 
    \[
    T_R(x) = \begin{cases}
        R & x > R \\
        x & -R \leq x \leq R \\
        -R & x < -R.
    \end{cases}
    \]
    It then follows that $|X_i - X_j| \geq |T_R(X_i) - T_R(X_j)|$, and therefore by the law of large numbers for U-statistics \citep{hoeffding1961strong}, we have
    \[
     \frac{1}{2n^2} \sum_{i,j=1}^n (X_i - X_j)^2 \geq \frac{1}{2n^2} \sum_{i,j=1}^n (T_R(X_i) - T_R(X_j))^2 \xrightarrow{\mathrm{a.s.}} \mathrm{Var}(T_R(X_1)).
    \]
    It thus suffices to prove that $V_\infty := \liminf_{R \to \infty} \mathrm{Var}(T_R(X_1)) = \infty$; to that end we proceed by contradiction. Suppose that $V_\infty  < \infty$. For any fixed $x \in \R$, pick $R > |x|$. Then if $|X_1| \leq x$, we have $|X_1| \leq R$ so $T_R(X_1) = X_1$ and thus 
    \begin{align}
        \P(|X_1| \leq x) &\leq  \P(|T_R(X_1)| \leq x) \notag \\ &\leq \P(|\E T_R(X_1)-T_R(X_1)| \geq \E T_R(X_1) - x) \notag\\ & \leq \frac{\mathrm{Var}(T_R(X_1))}{(x - \E T_R(X_1))^2}
        \label{eqn:X-trunc-ineq}
    \end{align}
    where the last inequality is Chebyshev's. Now since $T_R(X_1) \to X_1$ pointwise as $R \to \infty$, we have by Fatou's lemma that $\liminf_{R \to \infty} \E|T_R(X_1)| \geq \E |X_1| = \infty$.
    So pick any sequence $R_k \to \infty$ such that $\lim_{k \to \infty} \mathrm{Var}(T_{R_k}(X_1)) =V_\infty$ and take the limit $k \to \infty$ in (\ref{eqn:X-trunc-ineq}) to obtain $\P(X_1 \leq x) = 0$. But $x$ was arbitrary, so this is a contradiction. 
\end{proof}
Now we can prove Proposition~\ref{prop:divergence}.
\begin{proof}[Proof of Proposition~\ref{prop:divergence}]
     First suppose $M$ fails to exist at some $\lambda$, and $\E X < \infty$. Provided $n$ be large enough, $\lambda$ will be contained in $[-C_n, C_n]$, and by a law of large numbers for infinite means (\cite{durrett2019probability}, Theorem 2.3.8), $\tilde M_n(\lambda) \xrightarrow{\mathrm{a.s.}} \infty$, and then 
    \[
    \hat \xi_n^2 \geq 2\lambda^{-2}\psi_n(\lambda)  =  2\lambda^{-2}\log \tilde M_n(\lambda)-2\lambda^{-1}\bar X\xrightarrow{\mathrm{a.s.}}\infty.
    \]
    If $M$ does not exist at $\lambda$ but $\E X=\infty$, then we have by Proposition~\ref{prop:std-divergence} that $\hat \xi_n^2 \geq \hat \sigma_n^2\xrightarrow{\mathrm{a.s.}}\infty$.
    On the other hand, if $M$ exists everywhere, then for any $m > 0$ there is $\lambda$ such that $M(\lambda) > \exp(2^{-1}\lambda^2m^2)$. But $M_n(\lambda) \xrightarrow{\mathrm{a.s.}} M(\lambda)$, so provided $n$ be large enough, $M_n(\lambda) > \exp(2^{-1}\lambda^2m^2) - \eps$ occurs with probability 1. Then we have 
    \[
    \hat \xi^2_n \geq 2\lambda^{-2}\log M_n(\lambda) > m^2,
    \]
    so by the upper continuity of measure we have
    \[
    \P\left(\liminf_{n \to \infty} \hat \xi^2_n =\infty\right) = \P\left(\bigcap_{m=1}^\infty \liminf_{n \to \infty} \hat \xi_n^2 \geq m^2\right) = \lim_{m \to \infty} \P\left(\liminf_{n \to \infty}\hat \xi_n^2 \geq m^2\right) = 1.
    \]
\end{proof}


\begin{proof}[Proof of Proposition~\ref{prop:divergence2}]
For any $\lambda \in [0, C_n]$ we have
    \[
    \hat \xi_n^2 \geq \frac{2}{\lambda^2}\log\left(\frac1n\sum_{i=1}^n e^{\lambda(X_i - \bar X)}\right) \geq \frac{2}{\lambda^2}\log \left(\frac1ne^{(X_{(n)} - \bar X)\lambda}\right) = \frac{2}{\lambda}\left\{(X_{(n)} - \bar X)- \frac{\log n}{\lambda}\right\},
    \]
    and on the other hand if $\lambda \in [-C_n, 0]$ then
    \[
    \hat \xi_n^2 \geq \frac{2}{\lambda^2}\log\left(\frac1n\sum_{i=1}^n e^{\lambda(X_i - \bar X)}\right) \geq \frac{2}{\lambda^2}\log \left(\frac1ne^{(X_{(1)} - \bar X)\lambda}\right) \ge \frac{2}{|\lambda|}\left\{|X_{(1)} - \bar X|- \frac{\log n}{|\lambda|}\right\}.
    \]
    Note that $\Delta_n = | X_{(n)} - \bar{X}| \vee |X_{(1)} - \bar{X}|$. Assume without loss of generality that $\Delta_n = | X_{(n)} - \bar{X}|$. 

    We define $\lambda_n = (X_{(n)} - \bar X)^{-1}(2 \log n)$ and see that $\lambda_n \leq C_n$ under the condition that $\Delta_n \geq \frac{2\log n}{C_n}$. Plugging $\lambda_n$ into the RHS of the first inequality, we have
    \[
    \hat \xi_n^2 \geq \frac{(X_{(n)} - \bar X)^2}{2 \log n} = \frac{\Delta_n^2}{2 \log n}.
    \]
\end{proof}
\subsection{Details of Section 5}
\label{sec:experiments-appendix}
The algorithms were retrieved and implemented as described in \cite{liu2025beyond}. The permutation test was implemented through a modification to the Empirical Pipeline \citep{levi2021domino, liu2025beyond}, and the modified code can be accessed at \texttt{\href{https://github.com/LiuJ0/AMI-Benchmark}{https://github.com/LiuJ0/AMI-Benchmark}}.


\end{document}
